\documentclass[11pt,a4paper]{report}

% Packages
\usepackage[utf8]{inputenc}
\usepackage[T1]{fontenc}
\usepackage{amsmath,amssymb,amsthm}
\usepackage{mathtools}
\usepackage{bm}
\usepackage{xcolor}
\usepackage{listings}
\usepackage{hyperref}
\usepackage[margin=1in]{geometry}
\usepackage{tocloft}

% Theorem environments
\theoremstyle{definition}
\newtheorem{definition}{Definition}[section]
\newtheorem{theorem}[definition]{Theorem}
\newtheorem{lemma}[definition]{Lemma}
\newtheorem{proposition}[definition]{Proposition}
\newtheorem{corollary}[definition]{Corollary}
\newtheorem{example}[definition]{Example}
\newtheorem{remark}[definition]{Remark}

% Custom commands for Jordan algebra notation
\newcommand{\jmul}{\circ}           % Jordan product
\newcommand{\jsq}[1]{#1^{2}}        % Jordan square
\newcommand{\jpow}[2]{#1^{#2}}      % Jordan power
\newcommand{\jone}{\mathbf{1}}      % Jordan identity
\newcommand{\R}{\mathbb{R}}
\newcommand{\N}{\mathbb{N}}
\newcommand{\C}{\mathbb{C}}
\newcommand{\defas}{\coloneqq}      % Definition
\newcommand{\Lop}[1]{L_{#1}}        % Left multiplication operator
\newcommand{\Uop}[1]{U_{#1}}        % Quadratic operator
\newcommand{\Peirce}[2]{P_{#2}(#1)} % Peirce space

% Listings setup for Lean 4 code
\definecolor{leanblue}{RGB}{0,102,204}
\definecolor{leangreen}{RGB}{0,128,0}
\definecolor{leanpurple}{RGB}{128,0,128}
\definecolor{leangray}{RGB}{128,128,128}

\lstdefinelanguage{Lean}{
  keywords={def, theorem, lemma, example, structure, class, instance,
            where, extends, deriving, variable, namespace, end, open,
            import, section, axiom, abbrev, noncomputable, inductive,
            private, protected, mutual, attribute, set_option,
            universe, macro, syntax, elab},
  keywordstyle=\color{leanblue}\bfseries,
  keywords=[2]{by, fun, match, with, if, then, else, do, return, let, have, show, calc, at, intro, apply, exact, rfl, simp, rw, ring, omega, linarith, decide, sorry, suffices, refine, obtain, rcases, constructor, use, ext, induction, cases, contradiction, absurd, exfalso},
  keywordstyle=[2]\color{leanpurple},
  keywords=[3]{Type, Prop, Sort, Nat, Int, Real, Bool, True, False, And, Or, Not, Exists, Forall, Set, Finset, List, Option, Sum, Prod},
  keywordstyle=[3]\color{leangreen},
  sensitive=true,
  morecomment=[l]{--},
  morecomment=[n]{/-}{-/},
  morestring=[b]",
  commentstyle=\color{leangray}\itshape,
  stringstyle=\color{red},
  showstringspaces=false,
  basicstyle=\ttfamily\small,
  breaklines=true,
  breakatwhitespace=true,
  tabsize=2,
  frame=single,
  framerule=0.5pt,
  xleftmargin=0.5cm,
  xrightmargin=0.5cm,
}

\lstset{language=Lean}

% Title
\title{Jordan Algebra Formalization in Lean 4\\[0.5em]\large A Mathematical Rendering}
\author{Automated LaTeX Generation}
\date{\today}

\begin{document}

\maketitle

\tableofcontents

\chapter{Introduction}

This document presents a LaTeX rendering of the Jordan algebra formalization
from the \texttt{af-tests} Lean 4 project. Each section corresponds to a Lean
source file, presenting the formal definitions and theorems in mathematical notation.

\textbf{Notation:}
\begin{itemize}
  \item $a \jmul b$ denotes the Jordan product of $a$ and $b$
  \item $\jsq{a}$ denotes $a \jmul a$ (Jordan square)
  \item $\Lop{a}$ denotes the left multiplication operator $x \mapsto a \jmul x$
  \item $\Uop{a}$ denotes the quadratic operator $\Uop{a}(x) = 2\Lop{a}^2(x) - \Lop{\jsq{a}}(x)$
  \item $\Peirce{e}{\lambda}$ denotes the Peirce $\lambda$-eigenspace for idempotent $e$
\end{itemize}

\chapter{Core Definitions}
\section{Basic}

\begin{definition}[JordanAlgebra]
A \textbf{Jordan algebra} over $\R$ is a type $J$ with an additive commutative group structure and an $\R$-module structure, equipped with:
\begin{itemize}
  \item A bilinear product $\jmul : J \times J \to J$
  \item An identity element $\jone \in J$
\end{itemize}
satisfying:
\begin{enumerate}
  \item (Commutativity) $\forall a, b \in J$: $a \jmul b = b \jmul a$
  \item (Jordan identity) $\forall a, b \in J$: $(a \jmul b) \jmul (a \jmul a) = a \jmul (b \jmul (a \jmul a))$
  \item (Left identity) $\forall a \in J$: $\jone \jmul a = a$
  \item (Left distributivity) $\forall a, b, c \in J$: $a \jmul (b + c) = (a \jmul b) + (a \jmul c)$
  \item (Scalar compatibility) $\forall r \in \R$, $\forall a, b \in J$: $(r \cdot a) \jmul b = r \cdot (a \jmul b)$
\end{enumerate}
\end{definition}

\subsection{Basic Properties}

\begin{lemma}[jmul\_jone]
For all $a \in J$: $a \jmul \jone = a$.
\end{lemma}

\begin{lemma}[add\_jmul]
For all $a, b, c \in J$: $(a + b) \jmul c = (a \jmul c) + (b \jmul c)$.
\end{lemma}

\begin{lemma}[smul\_jmul]
For all $r \in \R$ and $a, b \in J$: $a \jmul (r \cdot b) = r \cdot (a \jmul b)$.
\end{lemma}

\begin{lemma}[jmul\_zero]
For all $a \in J$: $a \jmul 0 = 0$.
\end{lemma}

\begin{lemma}[zero\_jmul]
For all $a \in J$: $0 \jmul a = 0$.
\end{lemma}

\begin{lemma}[jmul\_neg]
For all $a, b \in J$: $a \jmul (-b) = -(a \jmul b)$.
\end{lemma}

\begin{lemma}[neg\_jmul]
For all $a, b \in J$: $(-a) \jmul b = -(a \jmul b)$.
\end{lemma}

\begin{lemma}[jmul\_sub]
For all $a, b, c \in J$: $a \jmul (b - c) = (a \jmul b) - (a \jmul c)$.
\end{lemma}

\begin{lemma}[sub\_jmul]
For all $a, b, c \in J$: $(a - b) \jmul c = (a \jmul c) - (b \jmul c)$.
\end{lemma}

\subsection{Square and Powers}

\begin{definition}[jsq]
The \textbf{Jordan square} of an element $a \in J$ is:
\[ \jsq{a} \defas a \jmul a \]
\end{definition}

\begin{lemma}[jsq\_def]
For all $a \in J$: $\jsq{a} = a \jmul a$.
\end{lemma}

\begin{theorem}[jordan\_identity']
For all $a, b \in J$:
\[ (a \jmul b) \jmul \jsq{a} = a \jmul (b \jmul \jsq{a}) \]
\end{theorem}

\begin{definition}[jpow]
The \textbf{Jordan power} of an element $a \in J$ is defined recursively:
\begin{align*}
\jpow{a}{0} &\defas \jone \\
\jpow{a}{n+1} &\defas a \jmul \jpow{a}{n}
\end{align*}
\end{definition}

\begin{lemma}[jpow\_zero]
For all $a \in J$: $\jpow{a}{0} = \jone$.
\end{lemma}

\begin{lemma}[jpow\_succ]
For all $a \in J$ and $n \in \N$: $\jpow{a}{n+1} = a \jmul \jpow{a}{n}$.
\end{lemma}

\begin{lemma}[jpow\_one]
For all $a \in J$: $\jpow{a}{1} = a$.
\end{lemma}

\begin{lemma}[jpow\_two]
For all $a \in J$: $\jpow{a}{2} = \jsq{a}$.
\end{lemma}

\subsection{Idempotents}

\begin{definition}[IsIdempotent]
An element $e \in J$ is \textbf{idempotent} if:
\[ \jsq{e} = e \]
\end{definition}

\begin{lemma}[IsIdempotent.jsq\_eq\_self]
If $e \in J$ is idempotent, then $\jsq{e} = e$.
\end{lemma}

\begin{theorem}[isIdempotent\_zero]
The zero element $0 \in J$ is idempotent.
\end{theorem}

\begin{theorem}[isIdempotent\_jone]
The identity element $\jone \in J$ is idempotent.
\end{theorem}

\begin{theorem}[IsIdempotent.complement]
If $e \in J$ is idempotent, then $\jone - e$ is also idempotent.
\end{theorem}

\section{Product}

\subsection{Left Multiplication Operator}

\begin{definition}[L]
The \textbf{left multiplication operator} by $a \in J$ is the $\R$-linear map $\Lop{a} : J \to J$ defined by:
\[ \Lop{a}(b) \defas a \jmul b \]
with linearity properties:
\begin{itemize}
  \item $\Lop{a}(b + c) = a \jmul b + a \jmul c$
  \item $\Lop{a}(r \cdot b) = r \cdot (a \jmul b)$
\end{itemize}
\end{definition}

\begin{lemma}[L\_apply]
For all $a, b \in J$:
\[ \Lop{a}(b) = a \jmul b \]
\end{lemma}

\begin{theorem}[L\_one]
The left multiplication operator by the identity is the identity map:
\[ \Lop{\jone} = \mathrm{id} \]
\end{theorem}

\begin{lemma}[L\_comm]
For all $a \in J$:
\[ \Lop{a} = \Lop{a} \]
\end{lemma}

\begin{definition}[Lsq]
The \textbf{squaring map} $\mathsf{Lsq} : J \to J$ is defined by:
\[ \mathsf{Lsq}(a) \defas \jsq{a} \]
\end{definition}

\begin{lemma}[Lsq\_eq\_L\_apply]
For all $a \in J$:
\[ \mathsf{Lsq}(a) = \Lop{a}(a) \]
\end{lemma}

\subsection{Jordan Identity for Operators}

\begin{theorem}[L\_comm\_L\_sq]
For all $a \in J$, the operators $\Lop{a}$ and $\Lop{\jsq{a}}$ commute:
\[ \Lop{a} \circ \Lop{\jsq{a}} = \Lop{\jsq{a}} \circ \Lop{a} \]
This is the operator form of the Jordan identity.
\end{theorem}

\begin{corollary}[L\_commute\_L\_sq]
For all $a \in J$:
\[ \mathrm{Commute}(\Lop{a}, \Lop{\jsq{a}}) \]
\end{corollary}

\subsection{Power Properties}

\begin{lemma}[jpow\_add\_one]
For all $a \in J$ and $n \in \N$:
\[ \jpow{a}{n+1} = a \jmul \jpow{a}{n} \]
\end{lemma}

\begin{lemma}[L\_jpow\_succ]
For all $a \in J$ and $n \in \N$:
\[ \jpow{a}{n+1} = \Lop{a}(\jpow{a}{n}) \]
\end{lemma}

\begin{lemma}[jmul\_jpow]
For all $a \in J$ and $n \in \N$, powers associate:
\[ a \jmul \jpow{a}{n} = \jpow{a}{n+1} \]
\end{lemma}

\begin{lemma}[jsq\_eq\_L]
For all $a \in J$, the square can be expressed in terms of $\Lop{}$:
\[ \jsq{a} = \Lop{a}(a) \]
\end{lemma}

\section{IsCommJordan}

\subsection{Multiplication Instance}

\begin{definition}[instMul]
Define multiplication on $J$ using the Jordan product:
\[ a \cdot b \defas a \jmul b \]
\end{definition}

\begin{lemma}[mul\_eq\_jmul]
For all $a, b \in J$:
\[ a \cdot b = a \jmul b \]
\end{lemma}

\subsection{Ring Structure}

\begin{proposition}[instNonUnitalNonAssocRing]
$J$ is an instance of $\mathsf{NonUnitalNonAssocRing}$, with:
\begin{itemize}
  \item (Left distributivity) $a \cdot (b + c) = a \cdot b + a \cdot c$
  \item (Right distributivity) $(a + b) \cdot c = a \cdot c + b \cdot c$
  \item (Zero multiplication) $0 \cdot a = 0$ and $a \cdot 0 = 0$
\end{itemize}
\end{proposition}

\begin{proposition}[instNonUnitalNonAssocCommRing]
$J$ is an instance of $\mathsf{NonUnitalNonAssocCommRing}$, with:
\begin{itemize}
  \item (Commutativity) $a \cdot b = b \cdot a$
\end{itemize}
\end{proposition}

\subsection{IsCommJordan Instance}

\begin{proposition}[instIsCommJordan]
$J$ is an instance of $\mathsf{IsCommJordan}$. The Jordan identity $(a \jmul b) \jmul a^2 = a \jmul (b \jmul a^2)$ is exactly what $\mathsf{IsCommJordan}$ requires.
\end{proposition}

\subsection{The Linearized Jordan Identity}

\begin{theorem}[linearized\_jordan\_mathlib]
For all $a, b, c \in J$:
\[ 2 \cdot \bigl( [L_a, L_{b \cdot c}] + [L_b, L_{c \cdot a}] + [L_c, L_{a \cdot b}] \bigr) = 0 \]
where $[f, g]$ denotes the Lie bracket (commutator) on $\mathsf{AddMonoid.End}$.
\end{theorem}

\begin{theorem}[linearized\_jordan\_jmul]
For all $a, b, c \in J$:
\[ 2 \cdot \bigl( [L_a, L_{b \jmul c}] + [L_b, L_{c \jmul a}] + [L_c, L_{a \jmul b}] \bigr) = 0 \]
\end{theorem}

\subsection{Connection: AddMonoid.End.mulLeft to L Operator}

\begin{lemma}[mulLeft\_apply]
For all $a, x \in J$:
\[ \mathsf{mulLeft}(a)(x) = a \jmul x \]
\end{lemma}

\begin{lemma}[mulLeft\_eq\_L\_apply]
For all $a, x \in J$:
\[ \mathsf{mulLeft}(a)(x) = \Lop{a}(x) \]
\end{lemma}

\begin{lemma}[lie\_mulLeft\_apply]
For all $f, g, x \in J$:
\[ [\mathsf{mulLeft}(f), \mathsf{mulLeft}(g)](x) = f \jmul (g \jmul x) - g \jmul (f \jmul x) \]
\end{lemma}

\section{Quadratic}

\subsection{The Quadratic U Operator}

\begin{definition}[U]
Let $J$ be a Jordan algebra and let $a, x \in J$. The \textbf{quadratic U operator} is defined by:
\[ \Uop{a}(x) \defas 2 \cdot (a \jmul (a \jmul x)) - \jsq{a} \jmul x \]
\end{definition}

\begin{theorem}[U\_def]
For all $a, x \in J$:
\[ \Uop{a}(x) = 2 \cdot (a \jmul (a \jmul x)) - \jsq{a} \jmul x \]
\end{theorem}

\subsection{Basic Properties}

\begin{theorem}[U\_one]
For all $x \in J$:
\[ \Uop{\jone}(x) = x \]
\end{theorem}

\begin{theorem}[U\_zero\_left]
For all $x \in J$:
\[ \Uop{0}(x) = 0 \]
\end{theorem}

\begin{theorem}[U\_zero\_right]
For all $a \in J$:
\[ \Uop{a}(0) = 0 \]
\end{theorem}

\begin{theorem}[U\_smul]
For all $c \in \R$ and $a, x \in J$:
\[ \Uop{c \cdot a}(x) = c^2 \cdot \Uop{a}(x) \]
\end{theorem}

\subsection{Linearity in the Second Argument}

\begin{theorem}[U\_add\_right]
For all $a, x, y \in J$:
\[ \Uop{a}(x + y) = \Uop{a}(x) + \Uop{a}(y) \]
\end{theorem}

\begin{theorem}[U\_smul\_right]
For all $a, x \in J$ and $c \in \R$:
\[ \Uop{a}(c \cdot x) = c \cdot \Uop{a}(x) \]
\end{theorem}

\begin{definition}[U\_linear]
Let $a \in J$. Define the $\R$-linear map $\Uop{a} : J \to_{\R} J$ by:
\begin{itemize}
  \item $\mathsf{toFun}(x) \defas \Uop{a}(x)$
  \item $\mathsf{map\_add'}$: $\Uop{a}$ is additive (by U\_add\_right)
  \item $\mathsf{map\_smul'}$: $\Uop{a}$ respects scalar multiplication (by U\_smul\_right)
\end{itemize}
\end{definition}

\begin{theorem}[U\_linear\_apply]
For all $a, x \in J$:
\[ \mathsf{U\_linear}(a)(x) = \Uop{a}(x) \]
\end{theorem}

\subsection{Relationship with L Operator}

\begin{theorem}[U\_eq\_L]
For all $a, x \in J$:
\[ \Uop{a}(x) = 2 \cdot \Lop{a}(\Lop{a}(x)) - \Lop{\jsq{a}}(x) \]
That is, $\Uop{a} = 2\Lop{a}^2 - \Lop{\jsq{a}}$.
\end{theorem}

\subsection{U Operator Properties}

\begin{theorem}[U\_self]
For all $a \in J$:
\[ \Uop{a}(a) = a \jmul \jsq{a} \]
\end{theorem}

\begin{theorem}[U\_idempotent\_self]
Let $e \in J$ be an idempotent (i.e., $\jsq{e} = e$). Then:
\[ \Uop{e}(e) = e \]
\end{theorem}

\begin{theorem}[U\_idempotent\_comp]
Let $e \in J$ be an idempotent. Then $\Uop{e}$ is idempotent as an operator:
\[ \mathsf{U\_linear}(e) \circ \mathsf{U\_linear}(e) = \mathsf{U\_linear}(e) \]
\textit{(Proof requires the fundamental formula.)}
\end{theorem}

\begin{theorem}[U\_L\_comm]
For all $a, x \in J$:
\[ \Uop{a}(\Lop{a}(x)) = \Lop{a}(\Uop{a}(x)) \]
\end{theorem}


\chapter{Operator Identities}
\section{LinearizedJordan}

\subsection{Four-Variable Identity}

\begin{theorem}[four\_variable\_identity]
For all $a, b, c, d \in J$:
\[
a \jmul ((b \jmul c) \jmul d) + b \jmul ((c \jmul a) \jmul d) + c \jmul ((a \jmul b) \jmul d) = (b \jmul c) \jmul (a \jmul d) + (c \jmul a) \jmul (b \jmul d) + (a \jmul b) \jmul (c \jmul d)
\]
\end{theorem}

\begin{theorem}[operator\_formula\_apply]
For all $a, b, c, d \in J$:
\[
a \jmul ((b \jmul c) \jmul d) + b \jmul ((c \jmul a) \jmul d) + c \jmul ((a \jmul b) \jmul d) = (a \jmul (b \jmul c)) \jmul d + b \jmul (a \jmul (c \jmul d)) + c \jmul (a \jmul (b \jmul d))
\]
\end{theorem}

\begin{theorem}[operator\_formula]
For all $a, b, c \in J$:
\[
\Lop{a} \circ \Lop{b \jmul c} + \Lop{b} \circ \Lop{c \jmul a} + \Lop{c} \circ \Lop{a \jmul b} = \Lop{a \jmul (b \jmul c)} + \Lop{b} \circ \Lop{a} \circ \Lop{c} + \Lop{c} \circ \Lop{a} \circ \Lop{b}
\]
\end{theorem}

\subsection{Power Associativity}

\begin{theorem}[L\_L\_jsq\_comm]
For all $a \in J$, the operators $\Lop{a}$ and $\Lop{\jsq{a}}$ commute.
\end{theorem}

\subsection{The Fundamental Jordan Identity (Element Form)}

\begin{theorem}[fundamental\_jordan]
For all $a, b \in J$:
\[
(\jsq{a} \jmul b) \jmul a = \jsq{a} \jmul (b \jmul a)
\]
\end{theorem}

\begin{theorem}[fundamental\_jordan']
For all $a, b \in J$:
\[
a \jmul (\jsq{a} \jmul b) = \jsq{a} \jmul (a \jmul b)
\]
\end{theorem}

\begin{theorem}[fundamental\_jordan\_original]
For all $a, b \in J$:
\[
(a \jmul b) \jmul \jsq{a} = a \jmul (b \jmul \jsq{a})
\]
\end{theorem}

\subsection{Power Associativity (Operator Commutation)}

\begin{theorem}[L\_jpow\_comm\_L]
For all $a \in J$ and $n \in \N$, the operators $\Lop{a}$ and $\Lop{\jpow{a}{n}}$ commute.
\end{theorem}

\begin{theorem}[jpow\_add]
For all $a \in J$ and $m, n \in \N$:
\[
\jpow{a}{m} \jmul \jpow{a}{n} = \jpow{a}{m+n}
\]
\end{theorem}

\begin{theorem}[jpow\_assoc]
For all $a \in J$ and $m, n, p \in \N$:
\[
(\jpow{a}{m} \jmul \jpow{a}{n}) \jmul \jpow{a}{p} = \jpow{a}{m} \jmul (\jpow{a}{n} \jmul \jpow{a}{p})
\]
\end{theorem}

\section{FundamentalFormula}

\begin{lemma}[linearized\_jordan\_aux]
For all $a, b, c \in J$:
\begin{align*}
&(a \jmul (b \jmul c)) \jmul \jsq{a} + (b \jmul (a \jmul c)) \jmul \jsq{a} + (c \jmul (a \jmul b)) \jmul \jsq{a} \\
&= a \jmul ((b \jmul c) \jmul \jsq{a}) + b \jmul ((a \jmul c) \jmul \jsq{a}) + c \jmul ((a \jmul b) \jmul \jsq{a})
\end{align*}
\end{lemma}

\begin{lemma}[U\_U\_expand]
For all $a, b, x \in J$:
\[ \Uop{a}(\Uop{b}(x)) = 2 \cdot a \jmul (a \jmul (2 \cdot b \jmul (b \jmul x) - \jsq{b} \jmul x)) - \jsq{a} \jmul (2 \cdot b \jmul (b \jmul x) - \jsq{b} \jmul x) \]
\end{lemma}

\subsection{The Fundamental Formula}

\begin{theorem}[fundamental\_formula]
\textbf{(Fundamental Formula)}
For all $a, b \in J$:
\[ \Uop{\Uop{a}(b)} = \Uop{a} \circ \Uop{b} \circ \Uop{a} \]
That is, for all $x \in J$:
\[ \Uop{\Uop{a}(b)}(x) = \Uop{a}(\Uop{b}(\Uop{a}(x))) \]
\end{theorem}

\subsection{Corollaries of the Fundamental Formula}

\begin{theorem}[U\_jsq]
For all $a \in J$:
\[ \Uop{\jsq{a}} = \Uop{a}^2 \]
That is, for all $x \in J$:
\[ \Uop{\jsq{a}}(x) = \Uop{a}(\Uop{a}(x)) \]
\end{theorem}

\begin{theorem}[U\_idempotent\_comp']
Let $e \in J$ be an idempotent (i.e., $\jsq{e} = e$). Then for all $x \in J$:
\[ \Uop{e}(\Uop{e}(x)) = \Uop{e}(x) \]
\end{theorem}

\begin{theorem}[U\_jsq\_idempotent]
Let $e \in J$ be an idempotent. Then for all $x \in J$:
\[ \Uop{\jsq{e}}(x) = \Uop{e}(x) \]
\end{theorem}

\subsection{Additional Properties from Fundamental Formula}

\begin{theorem}[U\_jpow\_two]
For all $a \in J$ and all $x \in J$:
\[ \Uop{\jpow{a}{2}}(x) = \Uop{a}(\Uop{a}(x)) \]
\end{theorem}

\begin{theorem}[fundamental\_formula\_linear]
For all $a, b \in J$, as $\R$-linear maps:
\[ \Uop{\Uop{a}(b)} = \Uop{a} \circ \Uop{b} \circ \Uop{a} \]
\end{theorem}

\section{OperatorIdentities}

\subsection{Commutator Bracket}

\begin{definition}[opComm]
The \textbf{commutator} (Lie bracket) of two $\R$-linear operators $f, g : J \to J$ is:
\[ \opComm{f}{g} \defas f \circ g - g \circ f \]
\end{definition}

\begin{notation}
We write $\llbracket f, g \rrbracket$ for $\opComm{f}{g}$.
\end{notation}

\begin{lemma}[opComm\_apply]
For all $\R$-linear maps $f, g : J \to J$ and $x \in J$:
\[ \llbracket f, g \rrbracket(x) = f(g(x)) - g(f(x)) \]
\end{lemma}

\begin{lemma}[opComm\_skew]
For all $\R$-linear maps $f, g : J \to J$:
\[ \llbracket f, g \rrbracket = -\llbracket g, f \rrbracket \]
\end{lemma}

\begin{lemma}[opComm\_self]
For all $\R$-linear maps $f : J \to J$:
\[ \llbracket f, f \rrbracket = 0 \]
\end{lemma}

\subsection{Linearized Jordan Identity}

\begin{theorem}[linearized\_jordan\_operator]
For all $a, b, c \in J$:
\[ 2 \cdot \bigl( \llbracket \Lop{a}, \Lop{b \jmul c} \rrbracket + \llbracket \Lop{b}, \Lop{c \jmul a} \rrbracket + \llbracket \Lop{c}, \Lop{a \jmul b} \rrbracket \bigr) = 0 \]
\end{theorem}

\begin{theorem}[linearized\_on\_jsq]
For all $a, b, c \in J$:
\begin{align*}
2 \cdot \bigl( &a \jmul ((b \jmul c) \jmul \jsq{a}) - (b \jmul c) \jmul (a \jmul \jsq{a}) \\
+ &b \jmul ((c \jmul a) \jmul \jsq{a}) - (c \jmul a) \jmul (b \jmul \jsq{a}) \\
+ &c \jmul ((a \jmul b) \jmul \jsq{a}) - (a \jmul b) \jmul (c \jmul \jsq{a}) \bigr) = 0
\end{align*}
\end{theorem}

\begin{theorem}[linearized\_core]
For all $a, b, c \in J$:
\begin{align*}
&a \jmul ((b \jmul c) \jmul \jsq{a}) - (b \jmul c) \jmul (a \jmul \jsq{a}) \\
+ &b \jmul ((c \jmul a) \jmul \jsq{a}) - (c \jmul a) \jmul (b \jmul \jsq{a}) \\
+ &c \jmul ((a \jmul b) \jmul \jsq{a}) - (a \jmul b) \jmul (c \jmul \jsq{a}) = 0
\end{align*}
\end{theorem}

\begin{theorem}[linearized\_rearranged]
For all $a, b, c \in J$:
\begin{align*}
&a \jmul ((b \jmul c) \jmul \jsq{a}) + b \jmul ((c \jmul a) \jmul \jsq{a}) + c \jmul ((a \jmul b) \jmul \jsq{a}) \\
&= (b \jmul c) \jmul (a \jmul \jsq{a}) + (c \jmul a) \jmul (b \jmul \jsq{a}) + (a \jmul b) \jmul (c \jmul \jsq{a})
\end{align*}
\end{theorem}

\subsection{Operator Commutator with Squared Element}

\begin{theorem}[operator\_commutator\_jsq]
For all $a, b \in J$:
\[ \llbracket \Lop{\jsq{a}}, \Lop{b} \rrbracket = 2 \cdot \llbracket \Lop{a}, \Lop{a \jmul b} \rrbracket \]
\end{theorem}

\begin{theorem}[operator\_commutator\_jsq\_apply]
For all $a, b, x \in J$:
\[ \jsq{a} \jmul (b \jmul x) - b \jmul (\jsq{a} \jmul x) = 2 \cdot \bigl( a \jmul ((a \jmul b) \jmul x) - (a \jmul b) \jmul (a \jmul x) \bigr) \]
\end{theorem}

\subsection{Idempotent Operator Identities}

\begin{theorem}[opComm\_idempotent\_eigenspace]
Let $e \in J$ be idempotent. Let $a \in J$ with $e \jmul a = \lambda \cdot a$ for some $\lambda \in \R$. Then for all $x \in J$:
\[ \llbracket \Lop{e}, \Lop{a} \rrbracket(x) = e \jmul (a \jmul x) - a \jmul (e \jmul x) \]
\end{theorem}

\begin{theorem}[L\_e\_L\_a\_L\_e]
Let $e \in J$ be idempotent. For all $a \in J$:
\[ \Lop{e} \circ \Lop{a} \circ \Lop{e} = \Lop{e \jmul (a \jmul e)} + \frac{1}{2} \cdot \llbracket \Lop{e}, \llbracket \Lop{e}, \Lop{a} \rrbracket \rrbracket \]
\emph{(Currently unproven in the formalization.)}
\end{theorem}

\begin{theorem}[opComm\_double\_idempotent]
Let $e \in J$ be idempotent. For all $a \in J$:
\[ \llbracket \Lop{e}, \llbracket \Lop{e}, \Lop{a} \rrbracket \rrbracket = 2 \cdot (\Lop{e} \circ \Lop{a} \circ \Lop{e}) - 2 \cdot \Lop{e \jmul (a \jmul e)} \]
\emph{(Currently unproven in the formalization.)}
\end{theorem}

\subsection{Useful Commutator Properties}

\begin{theorem}[opComm\_on\_P1]
Let $e \in J$ be idempotent. Let $a, b \in J$ with $e \jmul a = a$ and $e \jmul b = b$ (i.e., $a, b \in \Peirce{e}{1}$). Then:
\[ \llbracket \Lop{e}, \Lop{a} \rrbracket(b) = e \jmul (a \jmul b) - a \jmul b \]
\end{theorem}

\begin{theorem}[opComm\_on\_P0]
Let $e \in J$ be idempotent. Let $a \in J$ with $e \jmul a = 0$ (i.e., $a \in \Peirce{e}{0}$). Then for all $x \in J$:
\[ \llbracket \Lop{e}, \Lop{a} \rrbracket(x) = e \jmul (a \jmul x) - a \jmul (e \jmul x) \]
\end{theorem}


\chapter{Formally Real Jordan Algebras}
\section{Formally Real: Definitions}

\begin{definition}[FormallyRealJordan]
A Jordan algebra $J$ is \textbf{formally real} if for all $n \in \N$ and all functions $a : \mathrm{Fin}(n) \to J$:
\[
\sum_{i} \jsq{a_i} = 0 \implies \forall i,\, a_i = 0
\]
That is, sums of squares vanish only when all summands vanish.
\end{definition}

\begin{theorem}[sq\_eq\_zero\_imp\_zero]
Let $J$ be a formally real Jordan algebra. If $a \in J$ satisfies $\jsq{a} = 0$, then $a = 0$.
\end{theorem}

\begin{theorem}[of\_sq\_eq\_zero]
Let $J$ be a Jordan algebra. If for all $a \in J$ we have $\jsq{a} = 0 \implies a = 0$, then $J$ is formally real.
\end{theorem}

\begin{theorem}[formallyReal\_iff\_sq\_eq\_zero\_imp\_zero]
A Jordan algebra $J$ is formally real if and only if for all $a \in J$:
\[
\jsq{a} = 0 \implies a = 0
\]
\end{theorem}

\begin{definition}[FormallyRealJordan']
A Jordan algebra $J$ satisfies the \textbf{FormallyRealJordan'} property if for all $a \in J$:
\[
\jsq{a} = 0 \implies a = 0
\]
\end{definition}

\subsection{No Nilpotent Elements}

\begin{theorem}[jpow\_eq\_zero\_of\_ge]
Let $J$ be a Jordan algebra and $a \in J$. If $\jpow{a}{n} = 0$ and $m \ge n$, then $\jpow{a}{m} = 0$.
\end{theorem}

\begin{theorem}[no\_nilpotent\_of\_formallyReal]
Let $J$ be a formally real Jordan algebra. If $a \in J$ and $n \ge 1$ satisfy $\jpow{a}{n} = 0$, then $a = 0$.

That is, formally real Jordan algebras have no nilpotent elements.
\end{theorem}

\section{Formally Real: Properties}

\begin{theorem}[jsq\_ne\_zero]
Let $J$ be a Jordan algebra with the formally real property.
For all $a \in J$: if $a \ne 0$, then $\jsq{a} \ne 0$.
\end{theorem}

\begin{theorem}[jsq\_eq\_zero\_iff]
Let $J$ be a Jordan algebra with the formally real property.
For all $a \in J$:
\[ \jsq{a} = 0 \iff a = 0 \]
\end{theorem}

\begin{theorem}[smul\_jone\_ne\_zero']
Let $J$ be a Jordan algebra.
If $\jone \ne 0$, then for all $n \in \N$ with $n \ne 0$:
\[ n \cdot \jone \ne 0 \]
This is the Jordan algebra analog of characteristic zero.
\end{theorem}

\begin{theorem}[jpow\_two\_eq\_zero]
Let $J$ be a Jordan algebra with the formally real property.
For all $a \in J$: if $\jpow{a}{2} = 0$, then $a = 0$.
This is the fundamental ``no nilpotent elements'' property.
\end{theorem}

\begin{theorem}[subalgebra\_sq\_eq\_zero]
Let $J$ be a Jordan algebra with the formally real property, and let $a \in J$.
If $\jsq{a} = 0$, then $a = 0$.
Elements of a subalgebra satisfy the formally real property.
\end{theorem}

\section{Formally Real: Square Roots}

\subsection{Square Root Predicate}

\begin{definition}[IsPositiveSqrt]
Let $a, b \in J$. We say $b$ is a \textbf{positive square root} of $a$ if
\[ \jsq{b} = a \quad \land \quad b \in \mathsf{PositiveElement}(J). \]
\end{definition}

\begin{definition}[HasPositiveSqrt]
An element $a \in J$ \textbf{has a positive square root} if
\[ \exists\, b \in J \text{ such that } \mathsf{IsPositiveSqrt}(a, b). \]
\end{definition}

\subsection{Basic Properties}

\begin{lemma}[isPositiveSqrt\_def]
For $a, b \in J$:
\[ \mathsf{IsPositiveSqrt}(a, b) \iff \jsq{b} = a \land \mathsf{PositiveElement}(b). \]
\end{lemma}

\begin{theorem}[PositiveElement.of\_hasPositiveSqrt]
If $a \in J$ has a positive square root, then $a$ is a positive element.
\end{theorem}

\begin{lemma}[isPositiveSqrt\_zero]
Zero has zero as its positive square root:
\[ \mathsf{IsPositiveSqrt}(0, 0). \]
\end{lemma}

\begin{lemma}[isPositiveSqrt\_jone]
If $\jone$ is a positive element, then $\jone$ is its own positive square root:
\[ \mathsf{PositiveElement}(\jone) \implies \mathsf{IsPositiveSqrt}(\jone, \jone). \]
\end{lemma}

\subsection{Uniqueness}

\begin{theorem}[isPositiveSqrt\_unique]
Let $J$ be a formally real Jordan algebra. Let $a, b, c \in J$ with
\begin{align*}
\mathsf{IsPositiveSqrt}(a, b), \quad &\mathsf{IsPositiveSqrt}(a, c), \\
b \jmul c = c \jmul b, \quad &\jsq{b} \jmul c = c \jmul \jsq{b}.
\end{align*}
Then $b = c$.
\end{theorem}

\subsection{Existence}

\begin{theorem}[HasPositiveSqrt.of\_positiveElement]
Let $a \in J$ be a positive element. Then $a$ has a positive square root:
\[ \mathsf{PositiveElement}(a) \implies \mathsf{HasPositiveSqrt}(a). \]
\end{theorem}

\section{Formally Real: Ordered Cone}

\subsection{Positive Elements}

\begin{definition}[PositiveElement]
Let $J$ be a Jordan algebra. An element $a \in J$ is \textbf{positive} if it is a sum of squares:
\[ \mathsf{PositiveElement}(a) \defas \exists n \in \N,\, \exists b : \mathrm{Fin}(n) \to J,\quad a = \sum_{i=0}^{n-1} \jsq{b_i} \]
\end{definition}

\begin{theorem}[positiveElement\_zero]
Zero is a positive element:
\[ \mathsf{PositiveElement}(0) \]
\end{theorem}

\begin{theorem}[positiveElement\_jsq]
For all $a \in J$, the square $\jsq{a}$ is a positive element:
\[ \mathsf{PositiveElement}(\jsq{a}) \]
\end{theorem}

\begin{theorem}[positiveElement\_add]
The sum of positive elements is positive. If $\mathsf{PositiveElement}(a)$ and $\mathsf{PositiveElement}(b)$, then:
\[ \mathsf{PositiveElement}(a + b) \]
\end{theorem}

\begin{theorem}[positiveElement\_smul]
Positive scalar multiples of positive elements are positive. If $\mathsf{PositiveElement}(a)$ and $r > 0$, then:
\[ \mathsf{PositiveElement}(r \cdot a) \]
\end{theorem}

\subsection{Positive Cone}

\begin{definition}[positiveCone]
The \textbf{positive cone} in a Jordan algebra $J$ is the convex cone over $\R$ consisting of sums of squares:
\[ \mathsf{positiveCone}(J) : \mathsf{ConvexCone}\,\R\,J \]
where the carrier is $\{ a \in J \mid \mathsf{PositiveElement}(a) \}$, with:
\begin{itemize}
  \item $\mathsf{smul\_mem'}$: for $c > 0$ and $x$ in the cone, $c \cdot x$ is in the cone
  \item $\mathsf{add\_mem'}$: for $x, y$ in the cone, $x + y$ is in the cone
\end{itemize}
\end{definition}

\begin{theorem}[mem\_positiveCone]
For $a \in J$:
\[ a \in \mathsf{positiveCone}(J) \iff \mathsf{PositiveElement}(a) \]
\end{theorem}

\begin{theorem}[positiveCone\_pointed]
The positive cone is pointed (contains zero):
\[ \mathsf{positiveCone}(J).\mathsf{Pointed} \]
\end{theorem}

\begin{theorem}[positiveCone\_salient]
In a formally real Jordan algebra, the positive cone is salient. That is, if $J$ is formally real, then:
\[ \mathsf{positiveCone}(J).\mathsf{Salient} \]
In other words, if $a \in \mathsf{positiveCone}(J)$, $a \ne 0$, and $-a \in \mathsf{positiveCone}(J)$, we derive a contradiction.
\end{theorem}

\section{Formally Real: Spectrum}

\subsection{Orthogonal Idempotents}

\begin{definition}[AreOrthogonal]
Two elements $e, f \in J$ are \textbf{orthogonal} (in the Jordan sense) if their product is zero:
\[ \mathsf{AreOrthogonal}(e, f) \defas e \jmul f = 0 \]
\end{definition}

\begin{lemma}[AreOrthogonal.symm]
For all $e, f \in J$: if $\mathsf{AreOrthogonal}(e, f)$, then $\mathsf{AreOrthogonal}(f, e)$.
\end{lemma}

\begin{theorem}[jone\_sub\_idempotent\_orthogonal]
Let $e \in J$ be an idempotent. Then $e$ and $\jone - e$ are orthogonal:
\[ \mathsf{IsIdempotent}(e) \implies \mathsf{AreOrthogonal}(e, \jone - e) \]
\end{theorem}

\subsection{Complete System of Orthogonal Idempotents}

\begin{definition}[PairwiseOrthogonal]
A family of elements $(e_i)_{i \in \mathrm{Fin}(n)}$ is \textbf{pairwise orthogonal} if:
\[ \mathsf{PairwiseOrthogonal}(e) \defas \forall i, j,\, i \ne j \implies \mathsf{AreOrthogonal}(e_i, e_j) \]
\end{definition}

\begin{definition}[CSOI]
A \textbf{Complete System of Orthogonal Idempotents (CSOI)} of rank $n$ in a Jordan algebra $J$ consists of:
\begin{itemize}
  \item $\mathsf{idem} : \mathrm{Fin}(n) \to J$ --- the idempotents
  \item $\mathsf{is\_idem} : \forall i,\, \mathsf{IsIdempotent}(\mathsf{idem}(i))$ --- each element is idempotent
  \item $\mathsf{orthog} : \mathsf{PairwiseOrthogonal}(\mathsf{idem})$ --- pairwise orthogonal
  \item $\mathsf{complete} : \sum_{i} \mathsf{idem}(i) = \jone$ --- sum to identity
\end{itemize}
\end{definition}

\subsection{Spectral Decomposition}

\begin{definition}[SpectralDecomp]
A \textbf{spectral decomposition} of an element $a \in J$ consists of:
\begin{itemize}
  \item $n : \N$ --- the number of eigenvalues
  \item $\mathsf{eigenvalues} : \mathrm{Fin}(n) \to \R$ --- the eigenvalues
  \item $\mathsf{csoi} : \mathsf{CSOI}(J, n)$ --- the idempotent system
  \item $\mathsf{decomp} : a = \sum_{i} \mathsf{eigenvalues}(i) \cdot \mathsf{csoi}.\mathsf{idem}(i)$ --- decomposition equation
\end{itemize}
\end{definition}

\begin{definition}[jordanSpectrum]
The \textbf{spectrum} of an element $a \in J$ with spectral decomposition $\mathsf{sd}$ is:
\[ \mathsf{jordanSpectrum}(a, \mathsf{sd}) \defas \mathrm{range}(\mathsf{sd}.\mathsf{eigenvalues}) \subseteq \R \]
\end{definition}

\subsection{Properties of Orthogonal Idempotents}

\begin{theorem}[jmul\_orthog\_idem]
Let $e, f \in J$ be orthogonal elements, and let $r, s \in \R$. Then:
\[ (r \cdot e) \jmul (s \cdot f) = 0 \]
\end{theorem}

\begin{theorem}[jsq\_smul\_idem]
Let $e \in J$ be an idempotent and $r \in \R$. Then:
\[ \jsq{r \cdot e} = r^2 \cdot e \]
\end{theorem}

\subsection{Spectral Properties of Squares}

\begin{theorem}[spectral\_zero\_of\_eigenvalues\_zero]
Let $a \in J$ have spectral decomposition $\mathsf{sd}$. If all eigenvalues are zero:
\[ \bigl(\forall i,\, \mathsf{sd}.\mathsf{eigenvalues}(i) = 0\bigr) \implies a = 0 \]
\end{theorem}

\begin{theorem}[spectral\_eigenvalues\_sq\_nonneg]
For any spectral decomposition $\mathsf{sd}$ of $a \in J$:
\[ \forall i,\, 0 \le (\mathsf{sd}.\mathsf{eigenvalues}(i))^2 \]
\end{theorem}

\subsection{Formally Real and Spectral Properties}

\begin{theorem}[FormallyRealJordan.spectral\_sq\_eigenvalues\_nonneg]
Let $J$ be a formally real Jordan algebra and let $a \in J$. If $\jsq{a}$ has a spectral decomposition $\mathsf{sd}$, then all eigenvalues are non-negative:
\[ \forall i,\, 0 \le \mathsf{sd}.\mathsf{eigenvalues}(i) \]
\end{theorem}


\chapter{Peirce Decomposition}
\section{Peirce}

\subsection{Peirce Space Definition}

\begin{definition}[PeirceSpace]
Let $e \in J$ be an element of a Jordan algebra and $\lambda \in \R$. The \textbf{Peirce $\lambda$-eigenspace} for $e$ is:
\[ \Peirce{e}{\lambda} \defas \{ a \in J \mid e \jmul a = \lambda \cdot a \} \]
This is a submodule of $J$ over $\R$.
\end{definition}

\begin{lemma}[mem\_peirceSpace\_iff]
For all $e, a \in J$ and $\lambda \in \R$:
\[ a \in \Peirce{e}{\lambda} \iff e \jmul a = \lambda \cdot a \]
\end{lemma}

\begin{definition}[PeirceSpace$_0$, PeirceSpace$_{1/2}$, PeirceSpace$_1$]
The three standard Peirce spaces are:
\begin{itemize}
  \item $\Peirce{e}{0}$ --- the $0$-eigenspace
  \item $\Peirce{e}{1/2}$ --- the $\frac{1}{2}$-eigenspace
  \item $\Peirce{e}{1}$ --- the $1$-eigenspace
\end{itemize}
\end{definition}

\subsection{Basic Peirce Space Properties}

\begin{theorem}[peirceSpace\_zero\_eq\_ker]
For any $e \in J$:
\[ \Peirce{e}{0} = \ker(\Lop{e}) \]
\end{theorem}

\begin{theorem}[idempotent\_in\_peirce\_one]
If $e$ is idempotent (i.e., $\jsq{e} = e$), then $e \in \Peirce{e}{1}$.
\end{theorem}

\begin{theorem}[orthogonal\_in\_peirce\_zero]
If $e$ and $f$ are orthogonal (i.e., $e \jmul f = 0$), then $f \in \Peirce{e}{0}$.
\end{theorem}

\begin{theorem}[peirceSpace\_disjoint]
For any $e \in J$ and distinct eigenvalues $\lambda_1 \ne \lambda_2$, the Peirce spaces are disjoint:
\[ \Peirce{e}{\lambda_1} \cap \Peirce{e}{\lambda_2} = \{0\} \]
\end{theorem}

\begin{theorem}[jone\_in\_peirce\_one]
The identity element $\jone$ is in its own Peirce $1$-space:
\[ \jone \in \Peirce{\jone}{1} \]
\end{theorem}

\begin{theorem}[peirce\_one\_left\_id]
If $a \in \Peirce{e}{1}$, then $e$ acts as a left identity:
\[ e \jmul a = a \]
\end{theorem}

\begin{theorem}[peirce\_one\_right\_id]
If $a \in \Peirce{e}{1}$, then $e$ acts as a right identity:
\[ a \jmul e = a \]
\end{theorem}

\begin{theorem}[complement\_in\_peirce\_zero]
If $e$ is idempotent, then $\jone - e \in \Peirce{e}{0}$.
\end{theorem}

\subsection{Helper Lemmas for Peirce Polynomial}

\begin{lemma}[jmul\_jmul\_e\_x\_e]
For any idempotent $e$ and $x \in J$:
\[ (e \jmul x) \jmul e = e \jmul (e \jmul x) \]
\end{lemma}

\begin{lemma}[jmul\_nsmul]
For any $a, b \in J$ and $n \in \N$:
\[ a \jmul (n \cdot b) = n \cdot (a \jmul b) \]
\end{lemma}

\subsection{Peirce Polynomial Identity}

\begin{theorem}[peirce\_polynomial\_identity]
For any idempotent $e$, the left multiplication operator $\Lop{e}$ satisfies the Peirce polynomial:
\[ \Lop{e} \circ \left(\Lop{e} - \frac{1}{2} \cdot \mathrm{id}\right) \circ \left(\Lop{e} - \mathrm{id}\right) = 0 \]
Equivalently: $2\Lop{e}^3 - 3\Lop{e}^2 + \Lop{e} = 0$.

This fundamental identity shows that $\Lop{e}$ has eigenvalues only in $\{0, \frac{1}{2}, 1\}$.
\end{theorem}

\subsection{Peirce Multiplication Rules}

\begin{theorem}[peirce\_mult\_P0\_P0]
If $e$ is idempotent and $a, b \in \Peirce{e}{0}$, then:
\[ a \jmul b \in \Peirce{e}{0} \]
\end{theorem}

\begin{theorem}[peirce\_mult\_P1\_P1]
If $e$ is idempotent and $a, b \in \Peirce{e}{1}$, then:
\[ a \jmul b \in \Peirce{e}{1} \]
\end{theorem}

\begin{theorem}[peirce\_mult\_P0\_P1]
If $e$ is idempotent, $a \in \Peirce{e}{0}$, and $b \in \Peirce{e}{1}$, then:
\[ a \jmul b = 0 \]
\end{theorem}

\begin{theorem}[peirce\_mult\_P12\_P12]
If $e$ is idempotent and $a, b \in \Peirce{e}{1/2}$, then:
\[ a \jmul b \in \Peirce{e}{0} \oplus \Peirce{e}{1} \]
\end{theorem}

\begin{theorem}[peirce\_mult\_P0\_P12]
If $e$ is idempotent, $a \in \Peirce{e}{0}$, and $b \in \Peirce{e}{1/2}$, then:
\[ a \jmul b \in \Peirce{e}{1/2} \]
\end{theorem}

\begin{theorem}[peirce\_mult\_P1\_P12]
If $e$ is idempotent, $a \in \Peirce{e}{1}$, and $b \in \Peirce{e}{1/2}$, then:
\[ a \jmul b \in \Peirce{e}{1/2} \]
\end{theorem}

\subsection{Peirce Projection Operators}

\begin{definition}[peirceProj$_0$]
The projection onto $\Peirce{e}{0}$ is the Lagrange interpolation polynomial for eigenvalue $0$:
\[ \pi_0 \defas 2\Lop{e}^2 - 3\Lop{e} + \mathrm{id} = 2(\Lop{e} - \tfrac{1}{2}\mathrm{id})(\Lop{e} - \mathrm{id}) \]
\end{definition}

\begin{definition}[peirceProj$_{1/2}$]
The projection onto $\Peirce{e}{1/2}$ is the Lagrange interpolation polynomial for eigenvalue $\frac{1}{2}$:
\[ \pi_{1/2} \defas -4\Lop{e}^2 + 4\Lop{e} = -4\Lop{e}(\Lop{e} - \mathrm{id}) \]
\end{definition}

\begin{definition}[peirceProj$_1$]
The projection onto $\Peirce{e}{1}$ is the Lagrange interpolation polynomial for eigenvalue $1$:
\[ \pi_1 \defas 2\Lop{e}^2 - \Lop{e} = 2\Lop{e}(\Lop{e} - \tfrac{1}{2}\mathrm{id}) \]
\end{definition}

\begin{theorem}[peirceProj\_sum]
The three Peirce projections sum to the identity:
\[ \pi_0 + \pi_{1/2} + \pi_1 = \mathrm{id} \]
\end{theorem}

\begin{theorem}[peirceProj$_0$\_mem]
For any idempotent $e$ and $x \in J$:
\[ \pi_0(x) \in \Peirce{e}{0} \]
\end{theorem}

\begin{theorem}[peirceProj$_{1/2}$\_mem]
For any idempotent $e$ and $x \in J$:
\[ \pi_{1/2}(x) \in \Peirce{e}{1/2} \]
\end{theorem}

\begin{theorem}[peirceProj$_1$\_mem]
For any idempotent $e$ and $x \in J$:
\[ \pi_1(x) \in \Peirce{e}{1} \]
\end{theorem}

\subsection{Peirce Decomposition Theorem}

\begin{theorem}[peirce\_decomposition]
For any idempotent $e$ and $a \in J$, there exist unique elements $a_0, a_{1/2}, a_1 \in J$ such that:
\begin{enumerate}
  \item $a_0 \in \Peirce{e}{0}$
  \item $a_{1/2} \in \Peirce{e}{1/2}$
  \item $a_1 \in \Peirce{e}{1}$
  \item $a = a_0 + a_{1/2} + a_1$
\end{enumerate}
Moreover, $a_\lambda = \pi_\lambda(a)$ for $\lambda \in \{0, \frac{1}{2}, 1\}$.
\end{theorem}

\begin{theorem}[peirceSpace\_iSup\_eq\_top]
For any idempotent $e$:
\[ \Peirce{e}{0} \oplus \Peirce{e}{1/2} \oplus \Peirce{e}{1} = J \]
\end{theorem}

\begin{theorem}[peirce\_direct\_sum]
For any idempotent $e$, the Peirce spaces form an internal direct sum:
\[ J = \Peirce{e}{0} \oplus \Peirce{e}{1/2} \oplus \Peirce{e}{1} \]
That is, the decomposition is both complete (spanning) and independent (unique).
\end{theorem}

\subsection{PeirceOne Algebraic Structure}

\begin{definition}[PeirceOne]
For an idempotent $e$, define $\mathsf{PeirceOne}(e) \defas \Peirce{e}{1}$ as a subtype of $J$.
\end{definition}

\begin{proposition}[peirceOneAddCommGroup]
$\mathsf{PeirceOne}(e)$ inherits an abelian group structure from the submodule $\Peirce{e}{1}$.
\end{proposition}

\begin{proposition}[peirceOneModule]
$\mathsf{PeirceOne}(e)$ inherits an $\R$-module structure from the submodule $\Peirce{e}{1}$.
\end{proposition}

\begin{definition}[peirceOneMul]
For an idempotent $e$, define multiplication on $\mathsf{PeirceOne}(e)$ by restricting the Jordan product:
\[ a \cdot b \defas a \jmul b \]
This is well-defined by $\mathsf{peirce\_mult\_P1\_P1}$.
\end{definition}

\begin{definition}[peirceOneOne]
For an idempotent $e$, the identity element of $\mathsf{PeirceOne}(e)$ is $e$ itself.
\end{definition}

\begin{lemma}[peirceOne\_one\_val]
For an idempotent $e$:
\[ 1_{\mathsf{PeirceOne}(e)} = e \]
\end{lemma}

\begin{lemma}[peirceOne\_mul\_val]
For an idempotent $e$ and $a, b \in \mathsf{PeirceOne}(e)$:
\[ (a \cdot b).\mathsf{val} = a.\mathsf{val} \jmul b.\mathsf{val} \]
\end{lemma}

\begin{theorem}[peirceOne\_one\_mul]
For an idempotent $e$ and $a \in \mathsf{PeirceOne}(e)$:
\[ 1 \cdot a = a \]
\end{theorem}

\begin{theorem}[peirceOne\_mul\_one]
For an idempotent $e$ and $a \in \mathsf{PeirceOne}(e)$:
\[ a \cdot 1 = a \]
\end{theorem}

\begin{theorem}[peirceOne\_left\_distrib]
For an idempotent $e$ and $a, b, c \in \mathsf{PeirceOne}(e)$:
\[ a \cdot (b + c) = a \cdot b + a \cdot c \]
\end{theorem}

\begin{theorem}[peirceOne\_right\_distrib]
For an idempotent $e$ and $a, b, c \in \mathsf{PeirceOne}(e)$:
\[ (a + b) \cdot c = a \cdot c + b \cdot c \]
\end{theorem}

\begin{theorem}[peirceOne\_mul\_comm]
For an idempotent $e$ and $a, b \in \mathsf{PeirceOne}(e)$:
\[ a \cdot b = b \cdot a \]
\end{theorem}

\begin{proposition}[peirceOneFiniteDimensional]
If $J$ is a finite-dimensional Jordan algebra, then $\mathsf{PeirceOne}(e)$ is finite-dimensional.
\end{proposition}

\section{Primitive}

\subsection{Primitive Idempotent Definition}

\begin{definition}[IsPrimitive]
An idempotent $e \in J$ is \textbf{primitive} if:
\begin{enumerate}
  \item $e$ is idempotent,
  \item $e \ne 0$, and
  \item for all idempotent $f$ with $e \jmul f = f$, either $f = 0$ or $f = e$.
\end{enumerate}
\end{definition}

\begin{lemma}[IsPrimitive.isIdempotent]
If $e$ is primitive, then $e$ is idempotent.
\end{lemma}

\begin{lemma}[IsPrimitive.ne\_zero]
If $e$ is primitive, then $e \ne 0$.
\end{lemma}

\begin{lemma}[IsPrimitive.sub\_eq\_zero]
Let $e$ be primitive and let $f$ be idempotent with $e \jmul f = f$. Then $f = 0$ or $f = e$.
\end{lemma}

\begin{theorem}[isPrimitive\_of\_minimal]
Let $e$ be idempotent with $e \ne 0$. If for all nonzero idempotent $f$ with $e \jmul f = f$ we have $f = e$, then $e$ is primitive.
\end{theorem}

\subsection{Properties of Primitive Idempotents}

\begin{theorem}[IsPrimitive.orthog\_of\_ne]
Let $e$ be primitive and let $f$ be idempotent with $e \jmul f = f$ and $f \ne e$. Then $f = 0$.
\end{theorem}

\begin{theorem}[IsPrimitive.jmul\_complement]
Let $e$ be primitive. Then $e \jmul (\jone - e) = 0$.
\end{theorem}

\subsection{Field Classification Results (H-O 2.2.6)}

\begin{theorem}[finite\_dim\_field\_over\_real]
Let $F$ be a finite-dimensional field over $\R$. Then either $F \cong_{\R} \R$ or $F \cong_{\R} \C$.
\end{theorem}

\begin{theorem}[complex\_not\_formally\_real]
$\C$ is not formally real: there exist $n \in \N$ and $a : \mathrm{Fin}(n) \to \C$ such that $\sum_i a_i^2 = 0$ but $a_i \ne 0$ for some $i$.
\end{theorem}

\begin{theorem}[formallyReal\_field\_is\_real]
Let $F$ be a finite-dimensional field over $\R$. If $F$ is formally real (i.e., $\sum_i a_i^2 = 0$ implies $a_i = 0$ for all $i$), then $F \cong_{\R} \R$.
\end{theorem}

\subsection{Artinian Ring Structure Theorem (H-O 2.9.3)}

\begin{definition}[artinian\_reduced\_is\_product\_of\_fields]
Let $R$ be an Artinian reduced commutative ring. Then
\[ R \cong \prod_{I \in \mathrm{MaxSpec}(R)} R/I \]
where each $R/I$ is a field.
\end{definition}

\begin{proposition}[artinian\_reduced\_factor\_field]
Let $R$ be an Artinian reduced commutative ring and let $I \in \mathrm{MaxSpec}(R)$. Then $R/I$ is a field.
\end{proposition}

\subsection{Strongly Connected Idempotents}

\begin{definition}[IsStronglyConnected]
Two idempotents $e, f \in J$ are \textbf{strongly connected} if:
\begin{enumerate}
  \item $e$ and $f$ are orthogonal, and
  \item there exists $v \in J$ such that $e \jmul v = \frac{1}{2} v$, $f \jmul v = \frac{1}{2} v$, and $\jsq{v} = e + f$.
\end{enumerate}
\end{definition}

\subsection{Structure of Orthogonal Primitives (H-O 2.9.4(iv))}

\begin{theorem}[jpow\_succ\_mem\_peirce\_one]
Let $e$ be idempotent and let $a \in \Peirce{e}{1}$. Then for all $n \in \N$, $\jpow{a}{n+1} \in \Peirce{e}{1}$.
\end{theorem}

\begin{theorem}[primitive\_idempotent\_in\_P1]
Let $e$ be primitive and let $f$ be idempotent with $f \in \Peirce{e}{1}$. Then $f = 0$ or $f = e$.
\end{theorem}

\begin{theorem}[peirce\_one\_sq\_nonpos\_imp\_zero]
Let $J$ be formally real, let $e$ be idempotent with $e \ne 0$, and let $x \in \Peirce{e}{1}$. If $c \le 0$ and $\jsq{x} = c \cdot e$, then $x = 0$.
\end{theorem}

\begin{theorem}[lagrange\_idempotent\_in\_peirce\_one]
Let $J$ be formally real, let $e$ be idempotent, and let $x \in \Peirce{e}{1}$. Suppose $\jsq{x} = r \cdot x + s \cdot e$ for some $r, s \in \R$ with discriminant $\Delta = r^2 + 4s > 0$. Let $\mu = \frac{r - \sqrt{\Delta}}{2}$ and define
\[ f = \frac{1}{\sqrt{\Delta}} (x - \mu \cdot e). \]
Then $f$ is idempotent and $f \in \Peirce{e}{1}$.
\end{theorem}

\begin{theorem}[primitive\_peirce\_one\_dim\_one]
Let $J$ be a finite-dimensional formally real Jordan algebra and let $e$ be primitive. Then
\[ \dim_\R \Peirce{e}{1} = 1. \]
\end{theorem}

\begin{theorem}[primitive\_peirce\_one\_scalar]
Let $J$ be a finite-dimensional formally real Jordan algebra and let $e$ be primitive. For every $x \in \Peirce{e}{1}$, there exists $r \in \R$ such that $x = r \cdot e$.
\end{theorem}

\begin{theorem}[orthogonal\_primitive\_peirce\_sq]
Let $J$ be a finite-dimensional formally real Jordan algebra. Let $e, f$ be primitive and orthogonal, and let $a \in J$ satisfy $e \jmul a = \frac{1}{2} a$ and $f \jmul a = \frac{1}{2} a$. Then there exists $\mu \ge 0$ such that $\jsq{a} = \mu (e + f)$.
\end{theorem}

\begin{theorem}[orthogonal\_primitive\_structure]
Let $J$ be a finite-dimensional formally real Jordan algebra. Let $e, f$ be primitive and orthogonal. Then either:
\begin{enumerate}
  \item for all $a \in J$ with $e \jmul a = \frac{1}{2} a$ and $f \jmul a = \frac{1}{2} a$, we have $a = 0$, or
  \item $e$ and $f$ are strongly connected.
\end{enumerate}
\end{theorem}

\subsection{Decomposition into Primitives}

\begin{theorem}[exists\_primitive\_decomp]
Let $J$ be a finite-dimensional formally real Jordan algebra. Let $e$ be a nonzero idempotent. Then there exist $k \in \N$ and primitive idempotents $p_1, \ldots, p_k$ such that:
\begin{enumerate}
  \item each $p_i$ is primitive,
  \item the $p_i$ are pairwise orthogonal, and
  \item $e = \sum_{i=1}^{k} p_i$.
\end{enumerate}
\end{theorem}

\begin{theorem}[csoi\_refine\_primitive]
Let $J$ be a finite-dimensional formally real Jordan algebra and let $(e_1, \ldots, e_n)$ be a complete system of orthogonal idempotents (CSOI). Then there exists $m \ge n$ and a primitive CSOI $(p_1, \ldots, p_m)$ refining it, with each $p_i$ primitive.
\end{theorem}

\section{Eigenspace}

\subsection{Eigenspace Definition}

\begin{definition}[eigenspace]
Let $J$ be a Jordan algebra, $a \in J$, and $\mu \in \R$. The \textbf{$\mu$-eigenspace} of left multiplication by $a$ is defined as:
\[ \mathsf{eigenspace}(a, \mu) \defas \ker(\Lop{a} - \mu \cdot \mathrm{id}) \]
This is a submodule of $J$ consisting of elements $v$ such that $a \jmul v = \mu \cdot v$.
\end{definition}

\begin{lemma}[mem\_eigenspace\_iff]
For $a, v \in J$ and $\mu \in \R$:
\[ v \in \mathsf{eigenspace}(a, \mu) \iff a \jmul v = \mu \cdot v \]
\end{lemma}

\begin{lemma}[eigenspace\_zero]
For $a \in J$:
\[ \mathsf{eigenspace}(a, 0) = \ker(\Lop{a}) \]
\end{lemma}

\begin{theorem}[eigenspace\_eq\_peirceSpace]
For $a \in J$ and $\mu \in \R$:
\[ \mathsf{eigenspace}(a, \mu) = \Peirce{a}{\mu} \]
\end{theorem}

\subsection{Eigenvalue and Spectrum Definitions}

\begin{definition}[IsEigenvalue]
Let $a \in J$ and $\mu \in \R$. We say $\mu$ \textbf{is an eigenvalue of} $a$ if the $\mu$-eigenspace of $\Lop{a}$ is nontrivial:
\[ \mathsf{IsEigenvalue}(a, \mu) \defas \mathsf{eigenspace}(a, \mu) \ne \bot \]
\end{definition}

\begin{lemma}[isEigenvalue\_iff]
For $a \in J$ and $\mu \in \R$:
\[ \mathsf{IsEigenvalue}(a, \mu) \iff \mathsf{eigenspace}(a, \mu) \ne \bot \]
\end{lemma}

\begin{lemma}[isEigenvalue\_iff\_exists\_eigenvector]
For $a \in J$ and $\mu \in \R$:
\[ \mathsf{IsEigenvalue}(a, \mu) \iff \exists v \ne 0,\, a \jmul v = \mu \cdot v \]
\end{lemma}

\begin{definition}[eigenvalueSet]
The \textbf{eigenvalue set} of $a \in J$ is the set of all eigenvalues of $\Lop{a}$:
\[ \mathsf{eigenvalueSet}(a) \defas \{ \mu \in \R \mid \mathsf{IsEigenvalue}(a, \mu) \} \]
\end{definition}

\begin{lemma}[mem\_eigenvalueSet\_iff]
For $a \in J$ and $\mu \in \R$:
\[ \mu \in \mathsf{eigenvalueSet}(a) \iff \mathsf{IsEigenvalue}(a, \mu) \]
\end{lemma}

\subsection{Basic Eigenvalue Properties}

\begin{theorem}[eigenspace\_jone\_one]
The eigenspace of the identity element $\jone$ at eigenvalue $1$ is the entire algebra:
\[ \mathsf{eigenspace}(\jone, 1) = \top \]
\end{theorem}

\begin{theorem}[isEigenvalue\_jone\_one]
Let $J$ be a nontrivial Jordan algebra. Then $1$ is an eigenvalue of $\jone$:
\[ \mathsf{IsEigenvalue}(\jone, 1) \]
\end{theorem}

\begin{theorem}[eigenvalueSet\_jone]
Let $J$ be a nontrivial Jordan algebra. The eigenvalue set of $\jone$ is exactly $\{1\}$:
\[ \mathsf{eigenvalueSet}(\jone) = \{1\} \]
\end{theorem}

\subsection{Connection to Peirce Decomposition}

\begin{theorem}[idempotent\_eigenvalues\_subset]
Let $e \in J$ be an idempotent (i.e., $\jsq{e} = e$). Then the eigenvalues of $e$ are restricted to $\{0, \tfrac{1}{2}, 1\}$:
\[ \mathsf{eigenvalueSet}(e) \subseteq \left\{0, \tfrac{1}{2}, 1\right\} \]
This is a consequence of the Peirce polynomial identity $\Lop{e} \circ (\Lop{e} - \tfrac{1}{2}\mathrm{id}) \circ (\Lop{e} - \mathrm{id}) = 0$.
\end{theorem}

\subsection{Eigenspace Orthogonality}

\begin{theorem}[eigenspace\_orthogonal]
Let $J$ be a Jordan algebra with trace $\tau$. Let $a \in J$ and let $r, s \in \R$ with $r \ne s$. If $v \in \mathsf{eigenspace}(a, r)$ and $w \in \mathsf{eigenspace}(a, s)$, then $v$ and $w$ are orthogonal with respect to the trace inner product:
\[ \langle v, w \rangle_\tau = 0 \]
\end{theorem}

\begin{corollary}[eigenspace\_traceInner\_zero]
Let $J$ be a Jordan algebra with trace. For any $a \in J$ and distinct eigenvalues $r \ne s$:
\[ \forall v \in \mathsf{eigenspace}(a, r),\, \forall w \in \mathsf{eigenspace}(a, s),\, \langle v, w \rangle_\tau = 0 \]
\end{corollary}

\subsection{Finiteness in Finite Dimensions}

\begin{theorem}[eigenvalueSet\_finite]
Let $J$ be a finite-dimensional Jordan algebra. For any $a \in J$, the eigenvalue set is finite:
\[ \mathsf{eigenvalueSet}(a) \text{ is finite} \]
\end{theorem}


\chapter{Structure Theory}
\section{Simple}

\begin{definition}[IsSimpleJordan]
A Jordan algebra $J$ is \textbf{simple} if:
\begin{itemize}
  \item $\mathsf{nontrivial}$: There exists $a \in J$ such that $a \ne 0$.
  \item $\mathsf{ideal\_eq\_bot\_or\_top}$: For every Jordan ideal $I$ of $J$, either $I = \bot$ or $I = \top$.
\end{itemize}
\end{definition}

\begin{proposition}
If $J$ is a simple Jordan algebra, then $J$ is nontrivial.
\end{proposition}

\begin{theorem}[jone\_ne\_zero]
In a simple Jordan algebra $J$, the identity element is nonzero: $\jone \ne 0$.
\end{theorem}

\begin{theorem}[ideal\_eq\_top\_of\_ne\_bot]
Let $J$ be a simple Jordan algebra and let $I$ be a Jordan ideal of $J$. If $I \ne \bot$, then $I = \top$.
\end{theorem}

\begin{theorem}[ideal\_eq\_top\_of\_mem\_ne\_zero]
Let $J$ be a simple Jordan algebra, let $I$ be a Jordan ideal of $J$, and let $a \in J$. If $a \in I$ and $a \ne 0$, then $I = \top$.
\end{theorem}

\begin{theorem}[ideal\_of\_one\_eq\_top]
Let $J$ be a simple Jordan algebra and let $I$ be a Jordan ideal of $J$. If $\jone \in I$, then $I = \top$.
\end{theorem}

\begin{theorem}[IsSimpleJordan.mk']
Let $J$ be a Jordan algebra that is nontrivial. If every Jordan ideal $I$ of $J$ satisfies $I = \bot$ or $I = \top$, then $J$ is simple.
\end{theorem}

\begin{theorem}[isSimpleJordan\_of\_ideal\_trichotomy]
Let $J$ be a Jordan algebra. If there exists $a \in J$ with $a \ne 0$, and every Jordan ideal $I$ of $J$ satisfies $I = \bot$ or $I = \top$, then $J$ is simple.
\end{theorem}

\section{Semisimple}

\subsection{Ideal Operations}

\begin{definition}[idealSum]
Let $I, K$ be Jordan ideals of $J$. The \textbf{sum of two ideals} is defined as:
\[ \mathsf{idealSum}(I, K) \defas \{ x \in J \mid \exists\, a \in I,\, \exists\, b \in K,\, x = a + b \} \]
This forms a Jordan ideal with:
\begin{itemize}
  \item Addition: If $x = a_1 + b_1$ and $y = a_2 + b_2$ with $a_i \in I$, $b_i \in K$, then $x + y = (a_1 + a_2) + (b_1 + b_2)$
  \item Zero: $0 = 0 + 0$ with $0 \in I$ and $0 \in K$
  \item Scalar multiplication: $r \cdot x = r \cdot a + r \cdot b$
  \item Negation: $-x = (-a) + (-b)$
  \item Jordan multiplication: $y \jmul x = (y \jmul a) + (y \jmul b)$
\end{itemize}
\end{definition}

\begin{lemma}[mem\_idealSum]
For Jordan ideals $I, K$ and $x \in J$:
\[ x \in \mathsf{idealSum}(I, K) \iff \exists\, a \in I,\, \exists\, b \in K,\, x = a + b \]
\end{lemma}

\begin{lemma}[le\_idealSum\_left]
For Jordan ideals $I, K$: $I \subseteq \mathsf{idealSum}(I, K)$.
\end{lemma}

\begin{lemma}[le\_idealSum\_right]
For Jordan ideals $I, K$: $K \subseteq \mathsf{idealSum}(I, K)$.
\end{lemma}

\begin{definition}[idealInf]
Let $I, K$ be Jordan ideals of $J$. The \textbf{intersection of two ideals} is defined as:
\[ \mathsf{idealInf}(I, K) \defas I \cap K \]
This forms a Jordan ideal.
\end{definition}

\begin{lemma}[mem\_idealInf]
For Jordan ideals $I, K$ and $x \in J$:
\[ x \in \mathsf{idealInf}(I, K) \iff x \in I \land x \in K \]
\end{lemma}

\begin{lemma}[idealInf\_le\_left]
For Jordan ideals $I, K$: $\mathsf{idealInf}(I, K) \subseteq I$.
\end{lemma}

\begin{lemma}[idealInf\_le\_right]
For Jordan ideals $I, K$: $\mathsf{idealInf}(I, K) \subseteq K$.
\end{lemma}

\begin{definition}[Independent]
Two Jordan ideals $I, K$ are \textbf{independent} if their intersection is zero:
\[ \mathsf{Independent}(I, K) \defas \mathsf{idealInf}(I, K) = \bot \]
\end{definition}

\begin{lemma}[independent\_iff]
For Jordan ideals $I, K$:
\[ \mathsf{Independent}(I, K) \iff \forall\, x,\, (x \in I \to x \in K \to x = 0) \]
\end{lemma}

\begin{definition}[IsDirectSum]
Two Jordan ideals $I, K$ form a \textbf{direct sum decomposition} if they are independent and span $J$:
\[ \mathsf{IsDirectSum}(I, K) \defas \mathsf{Independent}(I, K) \land \mathsf{idealSum}(I, K) = \top \]
\end{definition}

\begin{lemma}[isDirectSum\_iff]
For Jordan ideals $I, K$:
\[ \mathsf{IsDirectSum}(I, K) \iff \bigl(\forall\, x,\, x \in I \to x \in K \to x = 0\bigr) \land \bigl(\forall\, y \in J,\, \exists\, a \in I,\, \exists\, b \in K,\, y = a + b\bigr) \]
\end{lemma}

\subsection{Semisimple Jordan Algebras}

\begin{definition}[IsSemisimpleJordan]
A Jordan algebra $J$ is \textbf{semisimple} if:
\begin{itemize}
  \item (Nontrivial) $\exists\, a \in J,\, a \ne 0$
  \item (Complement property) For every Jordan ideal $I$ of $J$, there exists a Jordan ideal $K$ such that $\mathsf{IsDirectSum}(I, K)$
\end{itemize}
\end{definition}

\begin{proposition}
Every semisimple Jordan algebra $J$ is nontrivial (i.e., $\exists\, a, b \in J$ with $a \ne b$).
\end{proposition}

\begin{theorem}[jone\_ne\_zero]
In a semisimple Jordan algebra $J$, the identity is nonzero: $\jone \ne 0$.
\end{theorem}
\begin{proof}
Suppose $\jone = 0$. By semisimplicity, there exists $a \in J$ with $a \ne 0$. But then:
\[ a = a \jmul \jone = a \jmul 0 = 0 \]
which contradicts $a \ne 0$.
\end{proof}

\begin{theorem}[IsSimpleJordan.isSemisimpleJordan]
Every simple Jordan algebra is semisimple.
\end{theorem}
\begin{proof}
Let $J$ be a simple Jordan algebra and $I$ a Jordan ideal. By simplicity, either $I = \bot$ or $I = \top$.
\begin{itemize}
  \item If $I = \bot$: Take $K = \top$. Then $\mathsf{idealInf}(\bot, \top) = \bot$ and $\mathsf{idealSum}(\bot, \top) = \top$.
  \item If $I = \top$: Take $K = \bot$. Then $\mathsf{idealInf}(\top, \bot) = \bot$ and $\mathsf{idealSum}(\top, \bot) = \top$.
\end{itemize}
In both cases, $\mathsf{IsDirectSum}(I, K)$ holds.
\end{proof}

\begin{theorem}[IsSemisimpleJordan.mk']
Let $J$ be a nontrivial Jordan algebra. If every Jordan ideal $I$ of $J$ has a complement $K$ such that $\mathsf{IsDirectSum}(I, K)$, then $J$ is semisimple.
\end{theorem}

\section{Ideal}

\begin{definition}[JordanIdeal]
A \textbf{Jordan ideal} of a Jordan algebra $J$ consists of:
\begin{itemize}
  \item $\mathsf{carrier} : \mathcal{P}(J)$ — a subset of $J$
  \item $\mathsf{add\_mem'}$ : $\forall a, b \in \mathsf{carrier},\ a + b \in \mathsf{carrier}$
  \item $\mathsf{zero\_mem'}$ : $0 \in \mathsf{carrier}$
  \item $\mathsf{smul\_mem'}$ : $\forall r \in \R,\ \forall a \in \mathsf{carrier},\ r \cdot a \in \mathsf{carrier}$
  \item $\mathsf{neg\_mem'}$ : $\forall a \in \mathsf{carrier},\ -a \in \mathsf{carrier}$
  \item $\mathsf{jmul\_mem'}$ : $\forall a \in J,\ \forall b \in \mathsf{carrier},\ a \jmul b \in \mathsf{carrier}$
\end{itemize}
\end{definition}

\begin{proposition}
$\mathsf{JordanIdeal}\ J$ is an instance of $\mathsf{SetLike}$.
\end{proposition}

\begin{lemma}[mem\_carrier]
For $I : \mathsf{JordanIdeal}\ J$ and $x : J$:
\[ x \in I.\mathsf{carrier} \iff x \in I \]
\end{lemma}

\begin{theorem}[add\_mem]
For $I : \mathsf{JordanIdeal}\ J$ and $a, b \in J$: if $a \in I$ and $b \in I$, then $a + b \in I$.
\end{theorem}

\begin{theorem}[zero\_mem]
For $I : \mathsf{JordanIdeal}\ J$: $0 \in I$.
\end{theorem}

\begin{theorem}[smul\_mem]
For $I : \mathsf{JordanIdeal}\ J$, $r \in \R$, and $a \in J$: if $a \in I$, then $r \cdot a \in I$.
\end{theorem}

\begin{theorem}[neg\_mem]
For $I : \mathsf{JordanIdeal}\ J$ and $a \in J$: if $a \in I$, then $-a \in I$.
\end{theorem}

\begin{theorem}[jmul\_mem]
For $I : \mathsf{JordanIdeal}\ J$, $a \in J$, and $b \in J$: if $b \in I$, then $a \jmul b \in I$.
\end{theorem}

\begin{theorem}[jmul\_mem\_left]
For $I : \mathsf{JordanIdeal}\ J$, $a \in J$, and $b \in J$: if $a \in I$, then $a \jmul b \in I$.
\end{theorem}

\begin{theorem}[sub\_mem]
For $I : \mathsf{JordanIdeal}\ J$ and $a, b \in J$: if $a \in I$ and $b \in I$, then $a - b \in I$.
\end{theorem}

\begin{definition}[toAddSubgroup]
For $I : \mathsf{JordanIdeal}\ J$, define $I.\mathsf{toAddSubgroup} : \mathsf{AddSubgroup}\ J$ with carrier $I.\mathsf{carrier}$, inheriting the additive group structure from $I$.
\end{definition}

\begin{definition}[toSubmodule]
For $I : \mathsf{JordanIdeal}\ J$, define $I.\mathsf{toSubmodule} : \mathsf{Submodule}\ \R\ J$ with carrier $I.\mathsf{carrier}$, inheriting the $\R$-module structure from $I$.
\end{definition}

\begin{proposition}
$\mathsf{JordanIdeal}\ J$ is an instance of $\mathsf{Bot}$.
\end{proposition}

\begin{definition}[bot]
The \textbf{zero ideal} $\bot : \mathsf{JordanIdeal}\ J$ is defined by:
\[ \mathsf{carrier} \defas \{0\} \]
\end{definition}

\begin{proposition}
$\mathsf{JordanIdeal}\ J$ is an instance of $\mathsf{Top}$.
\end{proposition}

\begin{definition}[top]
The \textbf{whole algebra ideal} $\top : \mathsf{JordanIdeal}\ J$ is defined by:
\[ \mathsf{carrier} \defas J \]
\end{definition}

\begin{lemma}[mem\_bot]
For $a \in J$: $a \in \bot \iff a = 0$.
\end{lemma}

\begin{lemma}[mem\_top]
For $a \in J$: $a \in \top$.
\end{lemma}

\section{Subalgebra}

\begin{definition}[JordanSubalgebra]
A \textbf{Jordan subalgebra} of a Jordan algebra $J$ consists of:
\begin{itemize}
  \item $\mathsf{carrier} : \mathcal{P}(J)$ --- a subset of $J$
  \item $\mathsf{jone\_mem'} : \jone \in \mathsf{carrier}$
  \item $\mathsf{jmul\_mem'} : \forall a, b,\; a \in \mathsf{carrier} \to b \in \mathsf{carrier} \to a \jmul b \in \mathsf{carrier}$
  \item $\mathsf{add\_mem'} : \forall a, b,\; a \in \mathsf{carrier} \to b \in \mathsf{carrier} \to a + b \in \mathsf{carrier}$
  \item $\mathsf{smul\_mem'} : \forall r \in \R,\; \forall a,\; a \in \mathsf{carrier} \to r \cdot a \in \mathsf{carrier}$
  \item $\mathsf{neg\_mem'} : \forall a,\; a \in \mathsf{carrier} \to -a \in \mathsf{carrier}$
\end{itemize}
\end{definition}

\begin{proposition}
$\mathsf{JordanSubalgebra}\;J$ is an instance of $\mathsf{SetLike}$.
\end{proposition}

\begin{lemma}[mem\_carrier]
For $S : \mathsf{JordanSubalgebra}\;J$ and $x : J$:
\[ x \in S.\mathsf{carrier} \iff x \in S \]
\end{lemma}

\begin{theorem}[jone\_mem]
For $S : \mathsf{JordanSubalgebra}\;J$:
\[ \jone \in S \]
\end{theorem}

\begin{theorem}[jmul\_mem]
For $S : \mathsf{JordanSubalgebra}\;J$ and $a, b : J$, if $a \in S$ and $b \in S$, then:
\[ a \jmul b \in S \]
\end{theorem}

\begin{theorem}[add\_mem]
For $S : \mathsf{JordanSubalgebra}\;J$ and $a, b : J$, if $a \in S$ and $b \in S$, then:
\[ a + b \in S \]
\end{theorem}

\begin{theorem}[smul\_mem]
For $S : \mathsf{JordanSubalgebra}\;J$, $r : \R$, and $a : J$, if $a \in S$, then:
\[ r \cdot a \in S \]
\end{theorem}

\begin{theorem}[neg\_mem]
For $S : \mathsf{JordanSubalgebra}\;J$ and $a : J$, if $a \in S$, then:
\[ -a \in S \]
\end{theorem}

\begin{theorem}[zero\_mem]
For $S : \mathsf{JordanSubalgebra}\;J$:
\[ 0 \in S \]
\end{theorem}

\begin{theorem}[sub\_mem]
For $S : \mathsf{JordanSubalgebra}\;J$ and $a, b : J$, if $a \in S$ and $b \in S$, then:
\[ a - b \in S \]
\end{theorem}

\begin{definition}[toAddSubgroup]
Let $S : \mathsf{JordanSubalgebra}\;J$. Define $\mathsf{toAddSubgroup}(S) : \mathsf{AddSubgroup}\;J$ by:
\begin{itemize}
  \item $\mathsf{carrier} \defas S.\mathsf{carrier}$
  \item $\mathsf{add\_mem'}$ --- inherited from $S$
  \item $\mathsf{zero\_mem'} \defas S.\mathsf{zero\_mem}$
  \item $\mathsf{neg\_mem'}$ --- inherited from $S$
\end{itemize}
\end{definition}

\begin{definition}[toSubmodule]
Let $S : \mathsf{JordanSubalgebra}\;J$. Define $\mathsf{toSubmodule}(S) : \mathsf{Submodule}\;\R\;J$ by:
\begin{itemize}
  \item $\mathsf{carrier} \defas S.\mathsf{carrier}$
  \item $\mathsf{add\_mem'}$ --- inherited from $S$
  \item $\mathsf{zero\_mem'} \defas S.\mathsf{zero\_mem}$
  \item $\mathsf{smul\_mem'}$ --- inherited from $S$
\end{itemize}
\end{definition}

\section{FiniteDimensional}

\subsection{Finite-Dimensional Class}

\begin{definition}[FinDimJordanAlgebra]
A \textbf{finite-dimensional Jordan algebra} is a Jordan algebra $J$ such that $J$ is finite-dimensional as an $\R$-vector space. Formally, it consists of:
\begin{itemize}
  \item $\mathsf{finDim} : \mathrm{FiniteDimensional}\ \R\ J$
\end{itemize}
\end{definition}

\subsection{Rank}

\begin{definition}[jordanRank]
The \textbf{Jordan rank} of a finite-dimensional Jordan algebra $J$ is:
\[ \mathsf{jordanRank}(J) \defas \mathrm{finrank}_\R(J) : \N \]
\end{definition}

\begin{theorem}[jordanRank\_pos]
Let $J$ be a finite-dimensional Jordan algebra. If $J$ is nontrivial, then:
\[ 0 < \mathsf{jordanRank}(J) \]
\end{theorem}

\subsection{Basis Existence}

\begin{theorem}[exists\_basis]
Let $J$ be a finite-dimensional Jordan algebra. Then there exists $n : \N$ and $b : \mathrm{Fin}(n) \to J$ such that $b$ is $\R$-linearly independent and $\mathrm{span}_\R(\mathrm{range}(b)) = \top$.
\end{theorem}

\begin{definition}[finBasis]
The \textbf{canonical basis} of a finite-dimensional Jordan algebra $J$ is:
\[ \mathsf{finBasis} : \mathrm{Basis}\ (\mathrm{Fin}(\mathsf{jordanRank}(J)))\ \R\ J \]
\end{definition}

\subsection{Linear Independence of Orthogonal Idempotents}

\begin{lemma}[jmul\_sum]
For any finite set $s : \mathrm{Finset}(\iota)$, element $e \in J$, scalars $r : \iota \to \R$, and elements $f : \iota \to J$:
\[ e \jmul \left( \sum_{i \in s} r_i \cdot f_i \right) = \sum_{i \in s} r_i \cdot (e \jmul f_i) \]
\end{lemma}

\begin{theorem}[linearIndependent\_orthog\_idem]
Let $n : \N$ and $e : \mathrm{Fin}(n) \to J$. If:
\begin{enumerate}
  \item $\forall i,\ e_i$ is idempotent,
  \item The family $(e_i)$ is pairwise orthogonal,
  \item $\forall i,\ e_i \ne 0$,
\end{enumerate}
then $(e_i)$ is $\R$-linearly independent.
\end{theorem}

\subsection{CSOI Cardinality Bound}

\begin{theorem}[csoi\_card\_le\_rank\_of\_nonzero]
Let $J$ be a finite-dimensional Jordan algebra and let $c$ be a CSOI of cardinality $n$. If $\forall i,\ c.\mathsf{idem}(i) \ne 0$, then:
\[ n \le \mathsf{jordanRank}(J) \]
\end{theorem}

\begin{theorem}[jone\_ne\_zero]
Let $J$ be a nontrivial Jordan algebra. Then:
\[ \jone \ne 0 \]
\end{theorem}

\begin{theorem}[csoi\_exists\_nonzero]
Let $J$ be a nontrivial Jordan algebra and let $c$ be a CSOI of cardinality $n$ with $0 < n$. Then:
\[ \exists i,\ c.\mathsf{idem}(i) \ne 0 \]
\end{theorem}

\section{Trace}

\subsection{Trace Class}

\begin{definition}[JordanTrace]
A \textbf{Jordan trace} on a Jordan algebra $J$ is a structure consisting of:
\begin{itemize}
  \item $\mathsf{trace} : J \to \R$ --- the trace functional
  \item $\mathsf{trace\_add}$: $\forall a, b \in J$, $\trace(a + b) = \trace(a) + \trace(b)$
  \item $\mathsf{trace\_smul}$: $\forall r \in \R$, $\forall a \in J$, $\trace(r \cdot a) = r \cdot \trace(a)$
  \item $\mathsf{trace\_jmul\_comm}$: $\forall a, b \in J$, $\trace(a \jmul b) = \trace(b \jmul a)$
  \item $\mathsf{trace\_L\_selfadjoint}$: $\forall a, v, w \in J$, $\trace((a \jmul v) \jmul w) = \trace(v \jmul (a \jmul w))$
  \item $\mathsf{trace\_jone\_pos}$: $0 < \trace(\jone)$
\end{itemize}
\end{definition}

\subsection{Basic Trace Properties}

\begin{theorem}[trace\_zero]
$\trace(0) = 0$.
\end{theorem}

\begin{theorem}[trace\_neg]
For all $a \in J$:
\[ \trace(-a) = -\trace(a) \]
\end{theorem}

\begin{theorem}[trace\_sub]
For all $a, b \in J$:
\[ \trace(a - b) = \trace(a) - \trace(b) \]
\end{theorem}

\subsection{Trace Inner Product}

\begin{definition}[traceInner]
The \textbf{trace inner product} is defined by:
\[ \tau(a, b) \defas \trace(a \jmul b) \]
\end{definition}

\begin{theorem}[traceInner\_symm]
For all $a, b \in J$:
\[ \tau(a, b) = \tau(b, a) \]
\end{theorem}

\begin{theorem}[traceInner\_add\_left]
For all $a, b, c \in J$:
\[ \tau(a + b, c) = \tau(a, c) + \tau(b, c) \]
\end{theorem}

\begin{theorem}[traceInner\_add\_right]
For all $a, b, c \in J$:
\[ \tau(a, b + c) = \tau(a, b) + \tau(a, c) \]
\end{theorem}

\begin{theorem}[traceInner\_smul\_left]
For all $r \in \R$ and $a, b \in J$:
\[ \tau(r \cdot a, b) = r \cdot \tau(a, b) \]
\end{theorem}

\begin{theorem}[traceInner\_smul\_right]
For all $r \in \R$ and $a, b \in J$:
\[ \tau(a, r \cdot b) = r \cdot \tau(a, b) \]
\end{theorem}

\begin{theorem}[traceInner\_zero\_left]
For all $a \in J$:
\[ \tau(0, a) = 0 \]
\end{theorem}

\begin{theorem}[traceInner\_zero\_right]
For all $a \in J$:
\[ \tau(a, 0) = 0 \]
\end{theorem}

\begin{theorem}[traceInner\_neg\_left]
For all $a, b \in J$:
\[ \tau(-a, b) = -\tau(a, b) \]
\end{theorem}

\begin{theorem}[traceInner\_neg\_right]
For all $a, b \in J$:
\[ \tau(a, -b) = -\tau(a, b) \]
\end{theorem}

\begin{theorem}[traceInner\_jmul\_left]
For all $a, v, w \in J$:
\[ \tau(a \jmul v, w) = \tau(v, a \jmul w) \]
That is, $\Lop{a}$ is self-adjoint with respect to the trace inner product.
\end{theorem}

\subsection{Non-degeneracy}

\begin{definition}[FormallyRealTrace]
A \textbf{formally real trace} on a Jordan algebra $J$ extends $\mathsf{JordanTrace}$ with:
\begin{itemize}
  \item $\mathsf{trace\_jsq\_nonneg}$: $\forall a \in J$, $0 \le \trace(\jsq{a})$
  \item $\mathsf{trace\_jsq\_eq\_zero\_iff}$: $\forall a \in J$, $\trace(\jsq{a}) = 0 \iff a = 0$
\end{itemize}
\end{definition}

\begin{theorem}[traceInner\_self\_nonneg]
For all $a \in J$:
\[ 0 \le \tau(a, a) \]
\end{theorem}

\begin{theorem}[traceInner\_self\_eq\_zero\_iff]
For all $a \in J$:
\[ \tau(a, a) = 0 \iff a = 0 \]
\end{theorem}

\begin{theorem}[traceInner\_self\_pos]
If $a \ne 0$, then:
\[ 0 < \tau(a, a) \]
\end{theorem}

\begin{theorem}[traceInner\_nondegenerate]
If $\tau(a, b) = 0$ for all $b \in J$, then $a = 0$.
\end{theorem}


\chapter{Spectral Theory}
\section{SpectralTheorem}

\subsection{Spectrum Definition}

\begin{definition}[spectrum]
The \textbf{spectrum} of an element $a \in J$ is the set of eigenvalues of $\Lop{a}$:
\[ \mathrm{spectrum}(a) \defas \mathrm{eigenvalueSet}(a) \]
In finite dimensions, this is a finite set of real numbers.
\end{definition}

\begin{lemma}[mem\_spectrum\_iff]
For $a \in J$ and $r \in \R$:
\[ r \in \mathrm{spectrum}(a) \iff \mathrm{IsEigenvalue}(a, r) \]
\end{lemma}

\begin{theorem}[spectrum\_finite]
Let $J$ be a finite-dimensional Jordan algebra. For any $a \in J$, $\mathrm{spectrum}(a)$ is finite.
\end{theorem}

\subsection{Spectral Decomposition Existence}

\begin{theorem}[spectral\_decomposition\_exists]
Let $J$ be a finite-dimensional formally real Jordan algebra with trace. For every $a \in J$, there exists a spectral decomposition $\mathsf{sd}$ of $a$ such that for all $i$, $\mathsf{sd}.\mathsf{csoi}.\mathsf{idem}(i)$ is primitive.
\end{theorem}

\begin{theorem}[spectral\_decomposition\_finset]
Let $J$ be a finite-dimensional formally real Jordan algebra with trace. For every $a \in J$, there exist:
\begin{itemize}
  \item A finite set $S \subseteq \R$
  \item A function $e : \R \to J$
\end{itemize}
such that:
\begin{enumerate}
  \item For all $r \in S$, $e(r)$ is idempotent
  \item For all $r, s \in S$ with $r \ne s$, $e(r)$ and $e(s)$ are orthogonal
  \item $\displaystyle\sum_{r \in S} e(r) = \jone$
  \item $\displaystyle a = \sum_{r \in S} r \cdot e(r)$
\end{enumerate}
\end{theorem}

\subsection{Uniqueness Results}

\begin{theorem}[spectrum\_eq\_eigenvalueSet]
Let $J$ be a finite-dimensional formally real Jordan algebra with trace. For $a \in J$ and any spectral decomposition $\mathsf{sd}$ of $a$:
\[ \mathrm{jordanSpectrum}(a, \mathsf{sd}) = \mathrm{spectrum}(a) \]
\end{theorem}

\begin{theorem}[spectral\_decomp\_eigenvalues\_eq\_spectrum]
Let $J$ be a finite-dimensional formally real Jordan algebra with trace. For $a \in J$ and spectral decomposition $\mathsf{sd}$:
\[ \mathrm{range}(\mathsf{sd}.\mathrm{eigenvalues}) = \mathrm{spectrum}(a) \]
\end{theorem}

\subsection{Square Eigenvalue Theorem}

\begin{lemma}[sum\_jmul]
For a finite index set $\iota$, elements $f : \iota \to J$, scalars $r : \iota \to \R$, and $e \in J$:
\[ \left(\sum_{i \in s} r_i \cdot f_i\right) \jmul e = \sum_{i \in s} r_i \cdot (f_i \jmul e) \]
\end{lemma}

\begin{theorem}[jsq\_sum\_orthog\_idem]
Let $c$ be a complete system of orthogonal idempotents of size $n$, and let $\mathrm{coef} : \mathrm{Fin}(n) \to \R$. Then:
\[ \jsq{\left(\sum_i \mathrm{coef}_i \cdot c.\mathrm{idem}(i)\right)} = \sum_i \mathrm{coef}_i^2 \cdot c.\mathrm{idem}(i) \]
\end{theorem}

\begin{theorem}[spectral\_sq]
Let $J$ be a finite-dimensional formally real Jordan algebra with trace. For $a \in J$ with spectral decomposition $\mathsf{sd}$, there exists a spectral decomposition $\mathsf{sd\_sq}$ of $\jsq{a}$ with the same size:
\[ \mathsf{sd\_sq}.n = \mathsf{sd}.n \]

The spectral decomposition of $\jsq{a}$ uses the same complete system of orthogonal idempotents with squared eigenvalues. If $a = \sum_i \lambda_i \cdot e_i$ with orthogonal idempotents, then:
\[ \jsq{a} = \left(\sum_i \lambda_i \cdot e_i\right)^2 = \sum_i \lambda_i^2 \cdot e_i \]
since $e_i \jmul e_j = 0$ for $i \ne j$ and $\jsq{e_i} = e_i$.
\end{theorem}

\begin{theorem}[spectrum\_sq]
Let $J$ be a finite-dimensional formally real Jordan algebra with trace. For $a \in J$:
\[ \mathrm{spectrum}(\jsq{a}) = \{r^2 : r \in \mathrm{spectrum}(a)\} \]
\end{theorem}

\subsection{Positivity of Square Eigenvalues}

\begin{theorem}[sq\_eigenvalues\_nonneg]
Let $J$ be a finite-dimensional formally real Jordan algebra with trace. For $a \in J$ and any spectral decomposition $\mathsf{sd}$ of $\jsq{a}$:
\[ \forall i,\; 0 \le \mathsf{sd}.\mathrm{eigenvalues}(i) \]
\end{theorem}

\begin{theorem}[spectrum\_sq\_nonneg]
Let $J$ be a finite-dimensional formally real Jordan algebra with trace. For $a \in J$:
\[ \forall r \in \mathrm{spectrum}(\jsq{a}),\; 0 \le r \]
\end{theorem}


\chapter{Reversibility}
\section{Reversible}

\begin{definition}[IsReversible]
A Jordan algebra $J$ is \textbf{reversible} if there exists:
\begin{itemize}
  \item An associative algebra $A$ over $\R$
  \item An $\R$-linear map $f : J \to A$
\end{itemize}
such that:
\begin{enumerate}
  \item (Jordan product preservation) For all $a, b \in J$:
    \[ f(a \jmul b) = \frac{1}{2}(f(a) \cdot f(b) + f(b) \cdot f(a)) \]
  \item (Identity preservation) $f(\jone) = 1_A$
  \item (Injectivity) $f$ is injective
  \item (Closure under reversal) For all $a, b, c \in J$, there exists $d \in J$ such that:
    \[ f(d) = f(a) \cdot f(b) \cdot f(c) + f(c) \cdot f(b) \cdot f(a) \]
\end{enumerate}
\end{definition}


\chapter{Classification}
\section{Classification: Simple Types}

\begin{definition}[SimpleType]
The \textbf{simple types} of finite-dimensional simple formally real Jordan algebras are given by the inductive type:
\begin{itemize}
  \item $\mathsf{typeI}(n) : \mathsf{SimpleType}$ --- $H_n(\R)$: real symmetric matrices
  \item $\mathsf{typeII}(n) : \mathsf{SimpleType}$ --- $H_n(\C)$: complex Hermitian matrices
  \item $\mathsf{typeIII}(n) : \mathsf{SimpleType}$ --- $H_n(\mathbb{H})$: quaternionic Hermitian matrices
  \item $\mathsf{typeIV}(n) : \mathsf{SimpleType}$ --- $V_n$: spin factors ($n \ge 2$)
  \item $\mathsf{typeV} : \mathsf{SimpleType}$ --- $H_3(\mathbb{O})$: Albert algebra (exceptional)
\end{itemize}
This type has decidable equality.
\end{definition}

\subsection{Properties of Simple Types}

\begin{definition}[dimension]
The \textbf{dimension} of a simple type is defined by:
\[
\mathsf{dimension}(t) \defas
\begin{cases}
\frac{n(n+1)}{2} & \text{if } t = \mathsf{typeI}(n) \\
n^2 & \text{if } t = \mathsf{typeII}(n) \\
n(2n-1) & \text{if } t = \mathsf{typeIII}(n) \\
n & \text{if } t = \mathsf{typeIV}(n) \\
27 & \text{if } t = \mathsf{typeV}
\end{cases}
\]
\end{definition}

\begin{definition}[isReversible]
A simple type is \textbf{reversible} (embeds into an associative algebra with reversal) if and only if it is not of type V:
\[
\mathsf{isReversible}(t) \defas
\begin{cases}
\mathsf{false} & \text{if } t = \mathsf{typeV} \\
\mathsf{true} & \text{otherwise}
\end{cases}
\]
\end{definition}

\begin{definition}[isExceptional]
A simple type is \textbf{exceptional} (not special) if and only if it is of type V:
\[
\mathsf{isExceptional}(t) \defas
\begin{cases}
\mathsf{true} & \text{if } t = \mathsf{typeV} \\
\mathsf{false} & \text{otherwise}
\end{cases}
\]
\end{definition}

\begin{definition}[minIndex]
The \textbf{minimum index} for which each type is defined:
\[
\mathsf{minIndex}(t) \defas
\begin{cases}
1 & \text{if } t = \mathsf{typeI}(n) \\
1 & \text{if } t = \mathsf{typeII}(n) \\
1 & \text{if } t = \mathsf{typeIII}(n) \\
2 & \text{if } t = \mathsf{typeIV}(n) \\
0 & \text{if } t = \mathsf{typeV}
\end{cases}
\]
Note: $V_1 \cong \R$ is not considered a ``spin'' factor, hence the minimum index of 2 for type IV.
\end{definition}

\section{Classification: Real Symmetric}

\begin{theorem}[nontrivial]
Let $n$ be a nonempty finite type. Then the Jordan algebra of real symmetric matrices $H_n(\R)$ is nontrivial:
\[ \exists A \in H_n(\R), \quad A \ne 0 \]
\end{theorem}

\begin{proposition}[isSimple]
Let $n$ be a nonempty finite type with decidable equality. Then the Jordan algebra of real symmetric matrices $H_n(\R)$ is a simple Jordan algebra:
\[ \mathsf{IsSimpleJordan}(H_n(\R)) \]
\end{proposition}

\section{Classification: Complex Hermitian}

\subsection{Nontriviality}

\begin{theorem}[nontrivial]
Let $n$ be a nonempty finite type. Then there exists a complex Hermitian matrix $A \in H_n(\C)$ such that $A \ne 0$.
\end{theorem}

\subsection{Simplicity}

\begin{proposition}[isSimple]
Let $n$ be a nonempty finite type with decidable equality. Then the Jordan algebra of complex Hermitian matrices $H_n(\C)$ is a simple Jordan algebra.
\end{proposition}


\chapter{Matrix Jordan Algebras}
\section{Matrix: Jordan Product}

\subsection{Basic Jordan Product}

Let $n$ be a finite type with decidable equality, and let $R$ be a field of characteristic zero.

\begin{definition}[jordanMul]
The \textbf{Jordan product} on matrices $A, B \in M_n(R)$ is defined by:
\[ A \circ B \defas \frac{1}{2}(AB + BA) \]
\end{definition}

\begin{lemma}[jordanMul\_def]
For all $A, B \in M_n(R)$:
\[ A \circ B = \frac{1}{2}(AB + BA) \]
\end{lemma}

\begin{theorem}[jordanMul\_comm]
The Jordan product is commutative. For all $A, B \in M_n(R)$:
\[ A \circ B = B \circ A \]
\end{theorem}

\begin{lemma}[jordanMul\_one]
The identity matrix is the identity for the Jordan product. For all $A \in M_n(R)$:
\[ A \circ I = A \]
\end{lemma}

\begin{lemma}[one\_jordanMul]
For all $A \in M_n(R)$:
\[ I \circ A = A \]
\end{lemma}

\begin{lemma}[jordanMul\_zero]
For all $A \in M_n(R)$:
\[ A \circ 0 = 0 \]
\end{lemma}

\begin{lemma}[zero\_jordanMul]
For all $A \in M_n(R)$:
\[ 0 \circ A = 0 \]
\end{lemma}

\begin{theorem}[jordanMul\_self]
The Jordan square equals the matrix square. For all $A \in M_n(R)$:
\[ A \circ A = A^2 \]
\end{theorem}

\subsection{Hermitian Matrices}

Let $R$ be a field of characteristic zero with a star ring structure.

\begin{theorem}[IsHermitian.jordanMul]
Hermitian matrices are closed under the Jordan product. If $A, B \in M_n(R)$ are Hermitian (i.e., $A^* = A$ and $B^* = B$), then $A \circ B$ is also Hermitian:
\[ (A \circ B)^* = A \circ B \]
\end{theorem}

\section{Matrix: Instance}

\begin{definition}[HermitianMatrix]
Let $n$ be a finite type and $R$ be an additive group with star structure. Define:
\[ \mathsf{HermitianMatrix}(n, R) \defas \mathsf{selfAdjoint}(\mathsf{Matrix}(n, n, R)) \]
i.e., the type of self-adjoint (Hermitian) $n \times n$ matrices over $R$.
\end{definition}

\subsection{Jordan Product on Hermitian Matrices}

\begin{definition}[jmul]
Let $A, B : \mathsf{HermitianMatrix}(n, R)$. Define the \textbf{Jordan product}:
\[ A \jmul B \defas \frac{1}{2}(AB + BA) \]
which is again Hermitian.
\end{definition}

\begin{lemma}[jmul\_val]
For $A, B : \mathsf{HermitianMatrix}(n, R)$:
\[ (A \jmul B).\mathsf{val} = \mathsf{jordanMul}(A.\mathsf{val}, B.\mathsf{val}) \]
\end{lemma}

\begin{theorem}[jmul\_comm]
The Jordan product on Hermitian matrices is commutative:
\[ A \jmul B = B \jmul A \]
\end{theorem}

\begin{definition}[one]
The \textbf{identity element} in $\mathsf{HermitianMatrix}(n, R)$ is:
\[ \jone \defas I_n \]
where $I_n$ is the identity matrix, which is Hermitian.
\end{definition}

\begin{lemma}[one\_val]
\[ \jone.\mathsf{val} = 1 \]
\end{lemma}

\begin{theorem}[jmul\_one]
For all $A : \mathsf{HermitianMatrix}(n, R)$:
\[ A \jmul \jone = A \]
\end{theorem}

\begin{theorem}[one\_jmul]
For all $A : \mathsf{HermitianMatrix}(n, R)$:
\[ \jone \jmul A = A \]
\end{theorem}

\begin{theorem}[jmul\_add]
The Jordan product distributes over addition:
\[ A \jmul (B + C) = (A \jmul B) + (A \jmul C) \]
\end{theorem}

\begin{theorem}[jmul\_smul]
The Jordan product respects scalar multiplication: for $r \in \R$,
\[ (r \cdot A) \jmul B = r \cdot (A \jmul B) \]
\end{theorem}

\subsection{Jordan Identity}

\begin{theorem}[jordanMul\_jordan\_identity]
For matrices $A, B : \mathsf{Matrix}(n, n, R)$, the Jordan identity holds:
\[ (A \jmul B) \jmul (A \jmul A) = A \jmul (B \jmul (A \jmul A)) \]
where $A \jmul B = \frac{1}{2}(AB + BA)$. Both sides expand to:
\[ \frac{1}{4}(ABA^2 + BA^3 + A^3B + A^2BA) \]
\end{theorem}

\begin{theorem}[jordan\_identity]
For Hermitian matrices $A, B : \mathsf{HermitianMatrix}(n, R)$:
\[ (A \jmul B) \jmul \jsq{A} = A \jmul (B \jmul \jsq{A}) \]
\end{theorem}

\subsection{Jordan Algebra Instance}

\begin{proposition}[jordanAlgebraHermitianMatrix]
Let $n$ be a finite decidable type and $R$ a field with star ring structure, characteristic zero, and $\R$-algebra structure with star module. Then $\mathsf{HermitianMatrix}(n, R)$ is a \textbf{Jordan algebra} with:
\begin{itemize}
  \item Product: $A \jmul B = \frac{1}{2}(AB + BA)$
  \item Commutativity: $A \jmul B = B \jmul A$
  \item Jordan identity: $(A \jmul B) \jmul \jsq{A} = A \jmul (B \jmul \jsq{A})$
  \item Identity: $\jone = I_n$
  \item Left identity: $\jone \jmul A = A$
  \item Distributivity: $A \jmul (B + C) = (A \jmul B) + (A \jmul C)$
  \item Scalar compatibility: $(r \cdot A) \jmul B = r \cdot (A \jmul B)$
\end{itemize}
\end{proposition}

\section{Matrix: Formally Real}

This section establishes that Hermitian matrices over an RCLike field form a formally real Jordan algebra.

\subsection{Jordan Square as Matrix Square}

\begin{theorem}[jsq\_val]
For a Hermitian matrix $A \in \mathsf{HermitianMatrix}(n, \mathbb{K})$, the Jordan square equals the matrix square:
\[ (\jsq{A}).\mathsf{val} = A.\mathsf{val} \cdot A.\mathsf{val} \]
\end{theorem}

\begin{theorem}[sq\_eq\_conjTranspose\_mul]
For a Hermitian matrix $A$, the matrix square equals the conjugate-transpose product:
\[ A.\mathsf{val} \cdot A.\mathsf{val} = A.\mathsf{val}^H \cdot A.\mathsf{val} \]
\end{theorem}

\begin{theorem}[sq\_posSemidef]
For a Hermitian matrix $A$, the matrix square is positive semidefinite:
\[ (A.\mathsf{val} \cdot A.\mathsf{val}).\mathsf{PosSemidef} \]
\end{theorem}

\begin{theorem}[jsq\_nonneg]
The Jordan square of a Hermitian matrix is nonnegative in the matrix order:
\[ 0 \le (\jsq{A}).\mathsf{val} \]
\end{theorem}

\subsection{Zero Characterization}

\begin{theorem}[sq\_eq\_zero\_imp\_zero]
For Hermitian matrices, if the square vanishes then the matrix is zero:
\[ A.\mathsf{val} \cdot A.\mathsf{val} = 0 \implies A = 0 \]
\end{theorem}

\begin{theorem}[jsq\_eq\_zero\_imp\_zero]
For Hermitian matrices, if the Jordan square vanishes then the matrix is zero:
\[ \jsq{A} = 0 \implies A = 0 \]
\end{theorem}

\subsection{Sum of Squares Characterization}

\begin{theorem}[sum\_jsq\_val]
For a family $a : \mathsf{Fin}(m) \to \mathsf{HermitianMatrix}(n, \mathbb{K})$, the sum of Jordan squares corresponds to the sum of matrix squares:
\[ \left(\sum_{i} \jsq{a_i}\right).\mathsf{val} = \sum_{i} a_i.\mathsf{val} \cdot a_i.\mathsf{val} \]
\end{theorem}

\begin{lemma}[each\_sq\_nonneg]
Each matrix square in a family is nonnegative:
\[ \forall i,\; 0 \le a_i.\mathsf{val} \cdot a_i.\mathsf{val} \]
\end{lemma}

\begin{theorem}[sum\_sq\_eq\_zero\_imp\_all\_zero]
If a sum of squares of Hermitian matrices is zero, each term is zero. This is the key property making Hermitian matrices formally real:
\[ \sum_{i} a_i.\mathsf{val} \cdot a_i.\mathsf{val} = 0 \implies \forall i,\; a_i = 0 \]
\end{theorem}

\subsection{FormallyRealJordan Instance}

\begin{proposition}[formallyRealJordanHermitianMatrix]
Hermitian matrices over an RCLike field form a formally real Jordan algebra:
\[ \mathsf{HermitianMatrix}(n, \mathbb{K}) \text{ is an instance of } \mathsf{FormallyRealJordan} \]
The instance satisfies the formally real property: for any finite family $a : \mathsf{Fin}(m) \to \mathsf{HermitianMatrix}(n, \mathbb{K})$,
\[ \sum_{i} \jsq{a_i} = 0 \implies \forall i,\; a_i = 0 \]
\end{proposition}

\section{Matrix: Trace}

\subsection{Trace Inner Product Definition}

\begin{definition}[traceInner]
The \textbf{trace inner product} on Hermitian matrices is defined as follows. Let $A, B : \mathsf{HermitianMatrix}(n, \mathbb{K})$. Define:
\[ \mathsf{traceInner}(A, B) \defas \Tr(A \cdot B) \]
\end{definition}

\begin{lemma}[traceInner\_def]
For all $A, B : \mathsf{HermitianMatrix}(n, \mathbb{K})$:
\[ \mathsf{traceInner}(A, B) = \Tr(A \cdot B) \]
\end{lemma}

\subsection{Basic Properties}

\begin{theorem}[traceInner\_comm]
The trace inner product is commutative. For all $A, B : \mathsf{HermitianMatrix}(n, \mathbb{K})$:
\[ \mathsf{traceInner}(A, B) = \mathsf{traceInner}(B, A) \]
\end{theorem}

\begin{theorem}[traceInner\_add\_left]
The trace inner product is additive in the first argument. For all $A, B, C : \mathsf{HermitianMatrix}(n, \mathbb{K})$:
\[ \mathsf{traceInner}(A + B, C) = \mathsf{traceInner}(A, C) + \mathsf{traceInner}(B, C) \]
\end{theorem}

\begin{theorem}[traceInner\_add\_right]
The trace inner product is additive in the second argument. For all $A, B, C : \mathsf{HermitianMatrix}(n, \mathbb{K})$:
\[ \mathsf{traceInner}(A, B + C) = \mathsf{traceInner}(A, B) + \mathsf{traceInner}(A, C) \]
\end{theorem}

\begin{theorem}[traceInner\_zero\_left]
For all $A : \mathsf{HermitianMatrix}(n, \mathbb{K})$:
\[ \mathsf{traceInner}(0, A) = 0 \]
\end{theorem}

\begin{theorem}[traceInner\_zero\_right]
For all $A : \mathsf{HermitianMatrix}(n, \mathbb{K})$:
\[ \mathsf{traceInner}(A, 0) = 0 \]
\end{theorem}

\subsection{Self Inner Product}

\begin{theorem}[traceInner\_self]
For all $A : \mathsf{HermitianMatrix}(n, \mathbb{K})$:
\[ \mathsf{traceInner}(A, A) = \Tr(A^2) \]
\end{theorem}

\begin{theorem}[sq\_eq\_conjTranspose\_mul']
For Hermitian $A$, we have $A^2 = A^\dagger A$. That is, for all $A : \mathsf{HermitianMatrix}(n, \mathbb{K})$:
\[ A \cdot A = A^\dagger \cdot A \]
\end{theorem}

\begin{theorem}[traceInner\_self\_conj]
The self inner product is real (equals its conjugate). For all $A : \mathsf{HermitianMatrix}(n, \mathbb{K})$:
\[ \overline{\mathsf{traceInner}(A, A)} = \mathsf{traceInner}(A, A) \]
\end{theorem}

\begin{theorem}[traceInner\_self\_nonneg]
The self inner product is nonnegative. For all $A : \mathsf{HermitianMatrix}(n, \mathbb{K})$:
\[ 0 \le \mathsf{traceInner}(A, A) \]
\end{theorem}

\begin{theorem}[traceInner\_self\_eq\_zero]
The self inner product is zero if and only if the matrix is zero. For all $A : \mathsf{HermitianMatrix}(n, \mathbb{K})$:
\[ \mathsf{traceInner}(A, A) = 0 \iff A = 0 \]
\end{theorem}

\subsection{Trace of Jordan Product}

\begin{theorem}[traceInner\_eq\_trace\_jmul]
The trace inner product equals the trace of the Jordan product. For all $A, B : \mathsf{HermitianMatrix}(n, \mathbb{K})$:
\[ \mathsf{traceInner}(A, B) = \Tr(A \jmul B) \]
where the Jordan product is $A \jmul B = \frac{1}{2}(AB + BA)$.
\end{theorem}

\begin{theorem}[traceInner\_conj]
The trace inner product is real-valued for all Hermitian pairs. For all $A, B : \mathsf{HermitianMatrix}(n, \mathbb{K})$:
\[ \overline{\mathsf{traceInner}(A, B)} = \mathsf{traceInner}(A, B) \]
\end{theorem}

\subsection{Real Inner Product}

\begin{definition}[traceInnerReal]
The \textbf{real-valued trace inner product} is defined as follows. Let $A, B : \mathsf{HermitianMatrix}(n, \mathbb{K})$. Define:
\[ \mathsf{traceInnerReal}(A, B) \defas \Re(\mathsf{traceInner}(A, B)) \]
where $\Re$ denotes the real part.
\end{definition}

\begin{lemma}[traceInnerReal\_def]
For all $A, B : \mathsf{HermitianMatrix}(n, \mathbb{K})$:
\[ \mathsf{traceInnerReal}(A, B) = \Re(\mathsf{traceInner}(A, B)) \]
\end{lemma}

\begin{theorem}[traceInnerReal\_comm]
The real trace inner product is commutative. For all $A, B : \mathsf{HermitianMatrix}(n, \mathbb{K})$:
\[ \mathsf{traceInnerReal}(A, B) = \mathsf{traceInnerReal}(B, A) \]
\end{theorem}

\begin{theorem}[traceInnerReal\_add\_left]
The real trace inner product is additive. For all $A, B, C : \mathsf{HermitianMatrix}(n, \mathbb{K})$:
\[ \mathsf{traceInnerReal}(A + B, C) = \mathsf{traceInnerReal}(A, C) + \mathsf{traceInnerReal}(B, C) \]
\end{theorem}

\begin{theorem}[traceInnerReal\_self\_nonneg]
The real trace inner product self is nonnegative. For all $A : \mathsf{HermitianMatrix}(n, \mathbb{K})$:
\[ 0 \le \mathsf{traceInnerReal}(A, A) \]
\end{theorem}

\begin{theorem}[traceInnerReal\_self\_eq\_zero]
The real trace inner product self equals zero if and only if $A = 0$. For all $A : \mathsf{HermitianMatrix}(n, \mathbb{K})$:
\[ \mathsf{traceInnerReal}(A, A) = 0 \iff A = 0 \]
\end{theorem}

\section{Matrix: Real Hermitian}

This section establishes the equivalence between Hermitian and symmetric matrices over $\R$, and defines the type of real symmetric matrices as a Jordan algebra.

\subsection{Hermitian and Symmetric Equivalence for Real Matrices}

\begin{theorem}[isHermitian\_iff\_isSymm]
Let $A : \mathrm{Matrix}(n, n, \R)$. Then $A$ is Hermitian if and only if $A$ is symmetric:
\[ A^* = A \iff A^\top = A \]
\end{theorem}

\begin{theorem}[conjTranspose\_eq\_transpose\_real]
For any real matrix $A : \mathrm{Matrix}(m, n, \R)$:
\[ A^* = A^\top \]
\end{theorem}

\begin{lemma}[IsSymm.isHermitian]
If $A : \mathrm{Matrix}(n, n, \R)$ is symmetric, then $A$ is Hermitian.
\end{lemma}

\begin{lemma}[IsHermitian.isSymm\_of\_real]
If $A : \mathrm{Matrix}(n, n, \R)$ is Hermitian, then $A$ is symmetric.
\end{lemma}

\subsection{Real Symmetric Matrix Type}

\begin{definition}[RealSymmetricMatrix]
Define $\mathsf{RealSymmetricMatrix}(n)$ as an abbreviation for $\mathsf{HermitianMatrix}(n, \R)$.

Over $\R$, this coincides with the set of matrices satisfying $A^\top = A$.
\end{definition}

\subsection{Basic Properties}

\begin{theorem}[transpose\_eq]
For any $A : \mathsf{RealSymmetricMatrix}(n)$:
\[ A^\top = A \]
\end{theorem}

\begin{definition}[ofSymm]
Let $A : \mathrm{Matrix}(n, n, \R)$ and let $h$ be a proof that $A$ is symmetric. Define:
\[ \mathsf{ofSymm}(A, h) : \mathsf{RealSymmetricMatrix}(n) \]
as the real symmetric matrix with underlying matrix $A$.
\end{definition}

\begin{lemma}[ofSymm\_val]
For any symmetric matrix $A : \mathrm{Matrix}(n, n, \R)$ with proof $h$:
\[ \mathsf{ofSymm}(A, h).\mathsf{val} = A \]
\end{lemma}

\begin{theorem}[isSymm]
For any $A : \mathsf{RealSymmetricMatrix}(n)$, the underlying matrix $A.\mathsf{val}$ is symmetric.
\end{theorem}

\subsection{Instances}

\begin{proposition}
$\mathsf{RealSymmetricMatrix}(n)$ is an instance of $\mathsf{JordanAlgebra}$.
\end{proposition}

\begin{proposition}
$\mathsf{RealSymmetricMatrix}(n)$ is an instance of $\mathsf{FormallyRealJordan}$.
\end{proposition}

\subsection{Special Matrices}

\begin{definition}[one]
The identity matrix in $\mathsf{RealSymmetricMatrix}(n)$:
\[ \mathsf{one} \defas \mathsf{HermitianMatrix.one} \]
\end{definition}

\begin{definition}[zero]
The zero matrix in $\mathsf{RealSymmetricMatrix}(n)$:
\[ \mathsf{zero} \defas 0 \]
\end{definition}

\begin{definition}[diagonal]
Let $d : n \to \R$ be a function. Define the diagonal matrix:
\[ \mathsf{diagonal}(d) : \mathsf{RealSymmetricMatrix}(n) \]
as the real symmetric matrix with diagonal entries given by $d$.
\end{definition}

\begin{lemma}[diagonal\_val]
For any $d : n \to \R$:
\[ \mathsf{diagonal}(d).\mathsf{val} = \mathrm{Matrix.diagonal}(d) \]
\end{lemma}

\subsection{Trace Inner Product}

\begin{definition}[traceInner]
For $A, B : \mathsf{RealSymmetricMatrix}(n)$, define the trace inner product:
\[ \langle A, B \rangle_{\mathrm{tr}} \defas \mathsf{HermitianMatrix.traceInnerReal}(A, B) \]
\end{definition}

\begin{theorem}[traceInner\_comm]
For all $A, B : \mathsf{RealSymmetricMatrix}(n)$:
\[ \langle A, B \rangle_{\mathrm{tr}} = \langle B, A \rangle_{\mathrm{tr}} \]
\end{theorem}

\begin{theorem}[traceInner\_self\_nonneg]
For all $A : \mathsf{RealSymmetricMatrix}(n)$:
\[ 0 \le \langle A, A \rangle_{\mathrm{tr}} \]
\end{theorem}

\section{Matrix: Complex Hermitian}

\begin{definition}[ComplexHermitianMatrix]
Let $n$ be a finite type. The \textbf{complex Hermitian matrices} are defined as:
\[ \mathsf{ComplexHermitianMatrix}(n) \defas \mathsf{HermitianMatrix}(n, \C) \]
These are matrices $A$ satisfying $A^\dagger = A$ where $A^\dagger$ denotes the conjugate transpose.
\end{definition}

\subsection{Basic Properties}

\begin{theorem}[conjTranspose\_eq]
Let $A : \mathsf{ComplexHermitianMatrix}(n)$. Then:
\[ A^\dagger = A \]
\end{theorem}

\begin{definition}[ofHermitian]
Let $A : \mathsf{Matrix}(n, n, \C)$ and $h : A^\dagger = A$. Define:
\[ \mathsf{ofHermitian}(A, h) : \mathsf{ComplexHermitianMatrix}(n) \defas \langle A, h \rangle \]
\end{definition}

\begin{lemma}[ofHermitian\_val]
For $A : \mathsf{Matrix}(n, n, \C)$ and $h : A^\dagger = A$:
\[ (\mathsf{ofHermitian}(A, h)).\mathsf{val} = A \]
\end{lemma}

\begin{theorem}[val\_isHermitian]
Let $A : \mathsf{ComplexHermitianMatrix}(n)$. Then $A.\mathsf{val}$ is Hermitian.
\end{theorem}

\subsection{Diagonal Entries are Real}

\begin{theorem}[diag\_re\_eq\_self]
Let $A : \mathsf{ComplexHermitianMatrix}(n)$ and $i : n$. Then:
\[ \Re(A_{ii}) = A_{ii} \]
That is, diagonal entries are real.
\end{theorem}

\begin{theorem}[diag\_re\_eq\_diag]
Let $A : \mathsf{ComplexHermitianMatrix}(n)$. Then:
\[ \lambda i.\, \Re(\mathsf{diag}(A)_i) = \mathsf{diag}(A) \]
\end{theorem}

\begin{theorem}[apply\_conj]
Let $A : \mathsf{ComplexHermitianMatrix}(n)$ and $i, j : n$. Then:
\[ \overline{A_{ji}} = A_{ij} \]
\end{theorem}

\subsection{Instances}

\begin{proposition}
$\mathsf{ComplexHermitianMatrix}(n)$ is an instance of $\mathsf{JordanAlgebra}$.
\end{proposition}

\begin{proposition}
$\mathsf{ComplexHermitianMatrix}(n)$ is an instance of $\mathsf{FormallyRealJordan}$.
\end{proposition}

\subsection{Special Matrices}

\begin{definition}[one]
The identity matrix in $\mathsf{ComplexHermitianMatrix}(n)$:
\[ \mathsf{one} \defas \mathsf{HermitianMatrix.one} \]
\end{definition}

\begin{definition}[zero]
The zero matrix in $\mathsf{ComplexHermitianMatrix}(n)$:
\[ \mathsf{zero} \defas \langle 0, \cdot \rangle \]
\end{definition}

\begin{definition}[realDiagonal]
Let $d : n \to \R$. The \textbf{real diagonal matrix} is:
\[ \mathsf{realDiagonal}(d) : \mathsf{ComplexHermitianMatrix}(n) \defas \mathsf{diagonal}(\lambda i.\, d_i) \]
where entries are embedded as complex numbers. This is Hermitian since real numbers are self-conjugate.
\end{definition}

\begin{lemma}[realDiagonal\_val]
For $d : n \to \R$:
\[ (\mathsf{realDiagonal}(d)).\mathsf{val} = \mathsf{diagonal}(\lambda i.\, d_i) \]
\end{lemma}

\begin{definition}[diagonal]
Let $d : n \to \C$ with $\forall i,\, \overline{d_i} = d_i$. The \textbf{diagonal matrix} is:
\[ \mathsf{diagonal}(d, hd) : \mathsf{ComplexHermitianMatrix}(n) \defas \langle \mathsf{Matrix.diagonal}(d), \cdot \rangle \]
\end{definition}

\subsection{Trace Inner Product}

\begin{definition}[traceInner]
The \textbf{trace inner product} for complex Hermitian matrices $A, B$:
\[ \mathsf{traceInner}(A, B) : \C \defas \mathsf{HermitianMatrix.traceInner}(A, B) \]
\end{definition}

\begin{definition}[traceInnerReal]
The \textbf{real trace inner product} for complex Hermitian matrices $A, B$:
\[ \mathsf{traceInnerReal}(A, B) : \R \defas \mathsf{HermitianMatrix.traceInnerReal}(A, B) \]
For Hermitian matrices, $\mathrm{Tr}(AB)$ is always real, so this extracts the real part.
\end{definition}

\begin{theorem}[traceInner\_comm]
For $A, B : \mathsf{ComplexHermitianMatrix}(n)$:
\[ \mathsf{traceInner}(A, B) = \mathsf{traceInner}(B, A) \]
\end{theorem}

\begin{theorem}[traceInnerReal\_comm]
For $A, B : \mathsf{ComplexHermitianMatrix}(n)$:
\[ \mathsf{traceInnerReal}(A, B) = \mathsf{traceInnerReal}(B, A) \]
\end{theorem}

\begin{theorem}[traceInnerReal\_self\_nonneg]
For $A : \mathsf{ComplexHermitianMatrix}(n)$:
\[ 0 \le \mathsf{traceInnerReal}(A, A) \]
\end{theorem}


\chapter{Quaternionic Jordan Algebras}
\section{Quaternion: Jordan Product}

\subsection{Jordan Product Definition}

\begin{definition}[jordanMul]
The \textbf{Jordan product} on quaternionic matrices is defined as follows. For matrices $A, B : \mathrm{Matrix}(n, n, \mathbb{H}_\R)$, define:
\[ \mathsf{jordanMul}(A, B) \defas 2^{-1} \cdot (A \cdot B + B \cdot A) \]
where $2^{-1}$ denotes real scalar multiplication.
\end{definition}

\begin{lemma}[jordanMul\_def]
For all $A, B : \mathrm{Matrix}(n, n, \mathbb{H}_\R)$:
\[ \mathsf{jordanMul}(A, B) = 2^{-1} \cdot (A \cdot B + B \cdot A) \]
\end{lemma}

\subsection{Jordan Product Properties}

\begin{theorem}[jordanMul\_comm]
The Jordan product is commutative. For all $A, B : \mathrm{Matrix}(n, n, \mathbb{H}_\R)$:
\[ \mathsf{jordanMul}(A, B) = \mathsf{jordanMul}(B, A) \]
\end{theorem}

\begin{theorem}[jordanMul\_one]
For all $A : \mathrm{Matrix}(n, n, \mathbb{H}_\R)$:
\[ \mathsf{jordanMul}(A, 1) = A \]
\end{theorem}

\begin{theorem}[one\_jordanMul]
For all $A : \mathrm{Matrix}(n, n, \mathbb{H}_\R)$:
\[ \mathsf{jordanMul}(1, A) = A \]
\end{theorem}

\begin{theorem}[jordanMul\_zero]
For all $A : \mathrm{Matrix}(n, n, \mathbb{H}_\R)$:
\[ \mathsf{jordanMul}(A, 0) = 0 \]
\end{theorem}

\begin{theorem}[zero\_jordanMul]
For all $A : \mathrm{Matrix}(n, n, \mathbb{H}_\R)$:
\[ \mathsf{jordanMul}(0, A) = 0 \]
\end{theorem}

\begin{theorem}[jordanMul\_self]
The Jordan square equals the matrix square. For all $A : \mathrm{Matrix}(n, n, \mathbb{H}_\R)$:
\[ \mathsf{jordanMul}(A, A) = A \cdot A \]
\end{theorem}

\subsection{Hermitian Preservation}

\begin{theorem}[hermitian\_jordanMul]
Quaternionic Hermitian matrices are closed under the Jordan product. If $A, B : \mathrm{Matrix}(n, n, \mathbb{H}_\R)$ are Hermitian, then $\mathsf{jordanMul}(A, B)$ is Hermitian.
\end{theorem}

\subsection{Jordan Product on QuaternionHermitianMatrix}

\begin{definition}[jmul]
The \textbf{Jordan product} on quaternionic Hermitian matrices is defined as follows. For $A, B : \mathsf{QuaternionHermitianMatrix}(n)$, define:
\[ \mathsf{jmul}(A, B) \defas \langle \mathsf{jordanMul}(A.\mathsf{val}, B.\mathsf{val}), \_ \rangle \]
where the proof witness follows from $\mathsf{hermitian\_jordanMul}$.
\end{definition}

\begin{lemma}[jmul\_val]
For all $A, B : \mathsf{QuaternionHermitianMatrix}(n)$:
\[ (\mathsf{jmul}(A, B)).\mathsf{val} = \mathsf{jordanMul}(A.\mathsf{val}, B.\mathsf{val}) \]
\end{lemma}

\begin{theorem}[jmul\_comm]
The Jordan product on $\mathsf{QuaternionHermitianMatrix}(n)$ is commutative. For all $A, B$:
\[ \mathsf{jmul}(A, B) = \mathsf{jmul}(B, A) \]
\end{theorem}

\begin{theorem}[jmul\_one]
For all $A : \mathsf{QuaternionHermitianMatrix}(n)$:
\[ \mathsf{jmul}(A, \mathsf{one}) = A \]
\end{theorem}

\begin{theorem}[one\_jmul]
For all $A : \mathsf{QuaternionHermitianMatrix}(n)$:
\[ \mathsf{jmul}(\mathsf{one}, A) = A \]
\end{theorem}

\section{Quaternion: Hermitian Matrices}

\subsection{CharZero Instance}

\begin{proposition}[Quaternion.instCharZero]
$\mathbb{H}_{\R}$ (quaternions over $\R$) has characteristic zero.
\end{proposition}

\subsection{StarModule Instance}

\begin{proposition}[Quaternion.instStarModule]
$\mathbb{H}_{\R}$ forms a $\mathsf{StarModule}$ over $\R$. That is, real scalar multiplication commutes with the star operation:
\[ \star(r \cdot q) = r \cdot \star(q) \]
for all $r \in \R$ and $q \in \mathbb{H}_{\R}$.
\end{proposition}

\subsection{Scalar-Star Commutativity}

\begin{theorem}[Quaternion.smul\_star\_comm]
For $r \in \R$ and $q \in \mathbb{H}_{\R}$:
\[ \star(r \cdot q) = r \cdot \star(q) \]
\end{theorem}

\subsection{Quaternionic Hermitian Matrix Type}

\begin{definition}[QuaternionHermitianMatrix]
For a finite type $n$, define
\[ \mathsf{QuaternionHermitianMatrix}(n) \defas \mathsf{HermitianMatrix}(n, \mathbb{H}_{\R}) \]
These are matrices $A$ satisfying $A^H = A$ where $(\cdot)^H$ denotes conjugate transpose using quaternion conjugation.
\end{definition}

\subsection{Basic Properties}

\begin{theorem}[conjTranspose\_eq]
Let $A$ be a quaternionic Hermitian matrix. Then $A^H = A$.
\end{theorem}

\begin{definition}[ofHermitian]
Let $A : \mathsf{Matrix}(n, n, \mathbb{H}_{\R})$ with $h : A^H = A$. Define
\[ \mathsf{ofHermitian}(A, h) : \mathsf{QuaternionHermitianMatrix}(n) \defas \langle A, h \rangle \]
\end{definition}

\begin{lemma}[ofHermitian\_val]
For $A : \mathsf{Matrix}(n, n, \mathbb{H}_{\R})$ with $h : A^H = A$:
\[ \mathsf{ofHermitian}(A, h).\mathsf{val} = A \]
\end{lemma}

\begin{theorem}[val\_isHermitian]
For $A : \mathsf{QuaternionHermitianMatrix}(n)$, the underlying matrix $A.\mathsf{val}$ is Hermitian.
\end{theorem}

\subsection{Diagonal Entries are Real}

\begin{theorem}[diag\_self\_conj]
Let $A : \mathsf{QuaternionHermitianMatrix}(n)$ and $i \in n$. Then
\[ \star(A_{ii}) = A_{ii} \]
That is, diagonal entries are self-conjugate (real quaternions).
\end{theorem}

\subsection{Conjugate Transpose Properties}

\begin{theorem}[conjTranspose\_smul\_real]
For $r \in \R$ and $A : \mathsf{Matrix}(n, n, \mathbb{H}_{\R})$:
\[ (r \cdot A)^H = r \cdot A^H \]
\end{theorem}

\begin{theorem}[conjTranspose\_mul\_hermitian]
Let $A, B : \mathsf{Matrix}(n, n, \mathbb{H}_{\R})$ be Hermitian. Then
\[ (A \cdot B)^H = B \cdot A \]
\end{theorem}

\subsection{Special Matrices}

\begin{definition}[one]
The identity matrix is a quaternionic Hermitian matrix:
\[ \mathsf{one} : \mathsf{QuaternionHermitianMatrix}(n) \defas \langle 1, \text{isHermitian\_one} \rangle \]
\end{definition}

\begin{definition}[zero]
The zero matrix is a quaternionic Hermitian matrix:
\[ \mathsf{zero} : \mathsf{QuaternionHermitianMatrix}(n) \defas \langle 0, \text{trivial} \rangle \]
\end{definition}

\begin{definition}[realDiagonal]
Let $d : n \to \R$. Define the real diagonal matrix:
\[ \mathsf{realDiagonal}(d) : \mathsf{QuaternionHermitianMatrix}(n) \defas \mathsf{diag}(d(i))_{i \in n} \]
where each $d(i) \in \R$ is embedded as a real quaternion.
\end{definition}

\begin{lemma}[realDiagonal\_val]
For $d : n \to \R$:
\[ \mathsf{realDiagonal}(d).\mathsf{val} = \mathsf{diag}((d(i) : \mathbb{H}_{\R}))_{i \in n} \]
\end{lemma}

\section{Quaternion: Instance}

\subsection{Additivity at Matrix Level}

\begin{theorem}[jordanMul\_add]
For all matrices $A, B, C \in M_n(\mathbb{H})$:
\[ A \bullet (B + C) = A \bullet B + A \bullet C \]
where $\bullet$ denotes the Jordan product.
\end{theorem}

\begin{theorem}[add\_jordanMul]
For all matrices $A, B, C \in M_n(\mathbb{H})$:
\[ (A + B) \bullet C = A \bullet C + B \bullet C \]
\end{theorem}

\subsection{Scalar Multiplication at Matrix Level}

\begin{theorem}[jordanMul\_smul]
For all $r \in \R$ and matrices $A, B \in M_n(\mathbb{H})$:
\[ (r \cdot A) \bullet B = r \cdot (A \bullet B) \]
\end{theorem}

\begin{theorem}[smul\_jordanMul]
For all $r \in \R$ and matrices $A, B \in M_n(\mathbb{H})$:
\[ A \bullet (r \cdot B) = r \cdot (A \bullet B) \]
\end{theorem}

\subsection{Jordan Identity at Matrix Level}

\begin{theorem}[jordanMul\_jordan\_identity]
For all matrices $A, B \in M_n(\mathbb{H})$:
\[ (A \bullet B) \bullet (A \bullet A) = A \bullet (B \bullet (A \bullet A)) \]
\end{theorem}

\subsection{Lifted Properties for Quaternionic Hermitian Matrices}

\begin{theorem}[jmul\_add]
For all quaternionic Hermitian matrices $A, B, C \in \mathsf{QuaternionHermitianMatrix}(n)$:
\[ A \jmul (B + C) = A \jmul B + A \jmul C \]
\end{theorem}

\begin{theorem}[add\_jmul]
For all quaternionic Hermitian matrices $A, B, C \in \mathsf{QuaternionHermitianMatrix}(n)$:
\[ (A + B) \jmul C = A \jmul C + B \jmul C \]
\end{theorem}

\begin{theorem}[jmul\_smul]
For all $r \in \R$ and quaternionic Hermitian matrices $A, B \in \mathsf{QuaternionHermitianMatrix}(n)$:
\[ (r \cdot A) \jmul B = r \cdot (A \jmul B) \]
\end{theorem}

\begin{theorem}[smul\_jmul]
For all $r \in \R$ and quaternionic Hermitian matrices $A, B \in \mathsf{QuaternionHermitianMatrix}(n)$:
\[ A \jmul (r \cdot B) = r \cdot (A \jmul B) \]
\end{theorem}

\subsection{Jordan Identity}

\begin{theorem}[jordan\_identity]
For all quaternionic Hermitian matrices $A, B \in \mathsf{QuaternionHermitianMatrix}(n)$:
\[ (A \jmul B) \jmul (A \jmul A) = A \jmul (B \jmul (A \jmul A)) \]
\end{theorem}

\subsection{JordanAlgebra Instance}

\begin{proposition}[moduleQuaternionHermitianMatrix]
$\mathsf{QuaternionHermitianMatrix}(n)$ is an instance of $\mathsf{Module}\ \R$.
\end{proposition}

\begin{proposition}[jordanAlgebraQuaternionHermitianMatrix]
$\mathsf{QuaternionHermitianMatrix}(n)$ is an instance of $\mathsf{JordanAlgebra}$, with:
\begin{itemize}
  \item $\mathsf{jmul} := \mathsf{QuaternionHermitianMatrix.jmul}$
  \item $\mathsf{jmul\_comm} := \mathsf{QuaternionHermitianMatrix.jmul\_comm}$
  \item $\mathsf{jordan\_identity} := \mathsf{QuaternionHermitianMatrix.jordan\_identity}$
  \item $\mathsf{jone} := \mathsf{QuaternionHermitianMatrix.one}$
  \item $\mathsf{jone\_jmul} := \mathsf{QuaternionHermitianMatrix.one\_jmul}$
  \item $\mathsf{jmul\_add} := \mathsf{QuaternionHermitianMatrix.jmul\_add}$
  \item $\mathsf{jmul\_smul} := \mathsf{QuaternionHermitianMatrix.jmul\_smul}$
\end{itemize}
\end{proposition}

\section{Quaternion: Formally Real}

\begin{theorem}[quaternion\_coe\_injective]
The coercion $(\cdot) : \R \to \mathbb{H}$ is injective.
\end{theorem}

\begin{lemma}[quaternion\_sum\_coe\_eq\_coe\_sum]
For any function $f : n \to \R$:
\[ \sum_{k \in n} (f(k) : \mathbb{H}) = \left(\sum_{k \in n} f(k)\right) : \mathbb{H} \]
\end{lemma}

\begin{theorem}[diag\_mul\_self\_hermitian]
Let $A : M_n(\mathbb{H})$ be a Hermitian matrix (i.e., $A^\dagger = A$). For any $i \in n$:
\[ (A \cdot A)_{ii} = \sum_{j \in n} \|A_{ij}\|^2 \]
where $\|\cdot\|^2$ denotes the quaternion norm squared.
\end{theorem}

\begin{theorem}[sum\_normSq\_eq\_zero\_iff]
Let $A : M_n(\mathbb{H})$ and $i \in n$. Then:
\[ \sum_{j \in n} \|A_{ij}\|^2 = 0 \iff \forall j,\, A_{ij} = 0 \]
\end{theorem}

\begin{theorem}[matrix\_sq\_eq\_zero\_imp\_zero]
Let $A : M_n(\mathbb{H})$ be Hermitian. If $A \cdot A = 0$, then $A = 0$.
\end{theorem}

\begin{theorem}[jsq\_val\_eq\_mul\_self]
For a quaternionic Hermitian matrix $A$, the Jordan square equals the matrix square:
\[ \jsq{A}.\mathsf{val} = A.\mathsf{val} \cdot A.\mathsf{val} \]
\end{theorem}

\begin{theorem}[sq\_eq\_zero\_imp\_zero]
Let $A$ be a quaternionic Hermitian matrix. If $\jsq{A} = 0$ in the Jordan algebra, then $A = 0$.
\end{theorem}

\begin{definition}[totalNormSq]
For a family of quaternionic Hermitian matrices $(A_k)_{k \in \mathrm{Fin}(m)}$ and diagonal index $i \in n$:
\[ \mathsf{totalNormSq}(A, i) \defas \sum_{k} \sum_{j} \|(A_k)_{ij}\|^2 \]
\end{definition}

\begin{lemma}[totalNormSq\_nonneg]
For any family $(A_k)$ and index $i$:
\[ 0 \le \mathsf{totalNormSq}(A, i) \]
\end{lemma}

\begin{theorem}[sum\_jsq\_diag]
For a family $(A_k)_{k \in \mathrm{Fin}(m)}$ of quaternionic Hermitian matrices and index $i \in n$:
\[ \left(\sum_k \jsq{A_k}\right)_{ii} = \sum_k (A_k \cdot A_k)_{ii} \]
\end{theorem}

\begin{theorem}[sum\_jsq\_diag\_eq\_totalNormSq]
For a family $(A_k)_{k \in \mathrm{Fin}(m)}$ of quaternionic Hermitian matrices and index $i \in n$:
\[ \left(\sum_k \jsq{A_k}\right)_{ii} = \mathsf{totalNormSq}(A, i) \]
\end{theorem}

\begin{theorem}[sum\_jsq\_zero\_imp\_totalNormSq\_zero]
Let $(A_k)_{k \in \mathrm{Fin}(m)}$ be a family of quaternionic Hermitian matrices. If $\sum_k \jsq{A_k} = 0$, then for all $i \in n$:
\[ \mathsf{totalNormSq}(A, i) = 0 \]
\end{theorem}

\begin{theorem}[sum\_sq\_eq\_zero\_imp\_all\_zero]
Let $(A_k)_{k \in \mathrm{Fin}(m)}$ be a family of quaternionic Hermitian matrices. If $\sum_k \jsq{A_k} = 0$, then $A_k = 0$ for all $k$.
\end{theorem}

\begin{proposition}[formallyReal]
The quaternionic Hermitian matrices $\mathsf{QuaternionHermitianMatrix}(n)$ form a formally real Jordan algebra.
\end{proposition}


\chapter{Spin Factors}
\section{Spin Factor: Definition}

\begin{definition}[SpinFactor]
The \textbf{spin factor} $V_n$ is defined as:
\[ \mathsf{SpinFactor}(n) \defas \R \times \R^n \]
where $\R^n$ is the $n$-dimensional Euclidean space.
\end{definition}

\begin{proposition}
$\mathsf{SpinFactor}(n)$ is an instance of $\mathsf{Inhabited}$, with default element $(0, 0)$.
\end{proposition}

\subsection{Jordan Product}

\begin{definition}[jordanMul]
The \textbf{Jordan product} on spin factors is defined by:
\[ (\alpha, v) \jmul (\beta, w) \defas (\alpha\beta + \langle v, w \rangle, \alpha w + \beta v) \]
where $\langle \cdot, \cdot \rangle$ denotes the inner product on $\R^n$.
\end{definition}

\begin{definition}[one]
The \textbf{identity element} of the spin factor is:
\[ \jone \defas (1, 0) \]
\end{definition}

\subsection{Component Accessors}

\begin{definition}[scalar]
The \textbf{scalar component} of a spin factor element $x = (\alpha, v)$ is:
\[ \mathsf{scalar}(x) \defas \alpha = x_1 \]
\end{definition}

\begin{definition}[vector]
The \textbf{vector component} of a spin factor element $x = (\alpha, v)$ is:
\[ \mathsf{vector}(x) \defas v = x_2 \]
\end{definition}

\subsection{Basic Properties}

\begin{theorem}[jordanMul\_comm]
The Jordan product on spin factors is commutative:
\[ \forall x, y \in V_n : x \jmul y = y \jmul x \]
\end{theorem}

\begin{theorem}[one\_jordanMul]
The identity element is a left identity:
\[ \forall x \in V_n : \jone \jmul x = x \]
\end{theorem}

\begin{theorem}[jordanMul\_one]
The identity element is a right identity:
\[ \forall x \in V_n : x \jmul \jone = x \]
\end{theorem}

\begin{theorem}[jordanMul\_add]
The Jordan product distributes over addition:
\[ \forall x, y, z \in V_n : x \jmul (y + z) = (x \jmul y) + (x \jmul z) \]
\end{theorem}

\begin{theorem}[jordanMul\_smul]
The Jordan product is compatible with scalar multiplication:
\[ \forall r \in \R, \forall x, y \in V_n : (r \cdot x) \jmul y = r \cdot (x \jmul y) \]
\end{theorem}

\subsection{Jordan Square}

\begin{definition}[sq]
The \textbf{square} of a spin factor element is:
\[ \mathsf{sq}(x) \defas x \jmul x \]
\end{definition}

\begin{lemma}[sq\_def]
For all $x \in V_n$:
\[ \mathsf{sq}(x) = x \jmul x \]
\end{lemma}

\begin{theorem}[sq\_scalar]
The scalar component of the square is:
\[ (\mathsf{sq}(x))_1 = x_1^2 + \langle x_2, x_2 \rangle \]
\end{theorem}

\begin{theorem}[sq\_vector]
The vector component of the square is:
\[ (\mathsf{sq}(x))_2 = (2 x_1) \cdot x_2 \]
\end{theorem}

\subsection{Jordan Identity}

\begin{theorem}[jordan\_identity]
The spin factor satisfies the Jordan identity:
\[ \forall x, y \in V_n : (x \jmul y) \jmul \mathsf{sq}(x) = x \jmul (y \jmul \mathsf{sq}(x)) \]
\end{theorem}

\subsection{Jordan Algebra Instance}

\begin{proposition}[jordanAlgebra]
$\mathsf{SpinFactor}(n)$ is an instance of $\mathsf{JordanAlgebra}$, with:
\begin{itemize}
  \item $\mathsf{jmul} \defas \mathsf{jordanMul}$
  \item $\mathsf{jmul\_comm} \defas \mathsf{jordanMul\_comm}$
  \item $\mathsf{jordan\_identity} \defas \mathsf{jordan\_identity}$
  \item $\mathsf{jone} \defas \mathsf{one}$
  \item $\mathsf{jone\_jmul} \defas \mathsf{one\_jordanMul}$
  \item $\mathsf{jmul\_add} \defas \mathsf{jordanMul\_add}$
  \item $\mathsf{jmul\_smul} \defas \mathsf{jordanMul\_smul}$
\end{itemize}
\end{proposition}

\begin{definition}[Notation]
The infix notation for the Jordan product on spin factors:
\[ x \circ_v y \defas \mathsf{jordanMul}(x, y) \]
\end{definition}

\section{Spin Factor: Formally Real}

\subsection{Auxiliary Lemmas}

\begin{lemma}[add\_eq\_zero\_of\_sq\_add\_nonneg]
Let $a, b \in \R$ with $a \ge 0$ and $b \ge 0$. If $a + b = 0$, then $a = 0$ and $b = 0$.
\end{lemma}

\subsection{Key Lemmas}

\begin{theorem}[sq\_eq\_zero\_imp\_scalar\_zero]
Let $x \in \mathsf{SpinFactor}_n$. If $\jsq{x} = 0$, then $x_1 = 0$.
\end{theorem}

\begin{theorem}[sq\_eq\_zero\_imp\_vector\_zero]
Let $x \in \mathsf{SpinFactor}_n$. If $\jsq{x} = 0$, then $x_2 = 0$.
\end{theorem}

\begin{theorem}[sq\_eq\_zero\_imp\_zero]
Let $x \in \mathsf{SpinFactor}_n$. If $\jsq{x} = 0$, then $x = 0$.
\end{theorem}

\subsection{Connection to Jordan Algebra}

\begin{lemma}[sq\_eq\_jsq]
For all $x \in \mathsf{SpinFactor}_n$:
\[ \mathsf{sq}(x) = \jsq{x} \]
\end{lemma}

\begin{theorem}[jsq\_scalar]
For $x \in \mathsf{SpinFactor}_n$, the scalar component of the Jordan square is:
\[ (\jsq{x})_1 = x_1^2 + \langle x_2, x_2 \rangle \]
\end{theorem}

\begin{theorem}[sum\_jsq\_scalar]
For a family $a : \mathsf{Fin}(m) \to \mathsf{SpinFactor}_n$:
\[ \left( \sum_{i} \jsq{a_i} \right)_1 = \sum_{i} \left( (a_i)_1^2 + \langle (a_i)_2, (a_i)_2 \rangle \right) \]
\end{theorem}

\begin{theorem}[jsq\_scalar\_nonneg]
For all $x \in \mathsf{SpinFactor}_n$:
\[ (\jsq{x})_1 \ge 0 \]
\end{theorem}

\begin{theorem}[sum\_jsq\_eq\_zero\_imp\_scalar\_zero]
Let $a : \mathsf{Fin}(m) \to \mathsf{SpinFactor}_n$. If $\sum_{i} \jsq{a_i} = 0$, then for all $i \in \mathsf{Fin}(m)$:
\[ (a_i)_1^2 + \langle (a_i)_2, (a_i)_2 \rangle = 0 \]
\end{theorem}

\begin{theorem}[scalar\_term\_zero\_imp\_zero]
Let $x \in \mathsf{SpinFactor}_n$. If $x_1^2 + \langle x_2, x_2 \rangle = 0$, then $x = 0$.
\end{theorem}

\subsection{Formally Real Instance}

\begin{proposition}[formallyReal]
$\mathsf{SpinFactor}_n$ is an instance of $\mathsf{FormallyRealJordan}$.
\end{proposition}


\end{document}
